\chapter[VGIT and Toric Birational Geometry]
{VGIT and Toric Birational\\ Geometry}
\label{chap:vgit-toric-birational}

\begin{chapteroverview}
The preceding chapters asked what a fan can tell us about a fixed toric
variety.  We now change the question: how can one toric variety change
birationally into another?

The Cox construction already contains the answer.  Fix an affine space
with a diagonal action of a torus \(G\).  In Chapter~6 we deleted a
fan-determined exceptional set and took a quotient.  Geometric invariant
theory gives a more flexible version of the same construction: a
character of \(G\) determines which points are allowed before we take the
quotient.  As the character moves, the allowed set changes only when we
cross certain walls.  The resulting chamber decomposition is the
\emph{secondary fan}.

Our goal in this final chapter is to make the progression
\[
\boxed{
\begin{gathered}
\text{Cox action and a character}
\longrightarrow
\text{GIT quotient}
\longrightarrow
\text{secondary fan}
\\[-1mm]
\longrightarrow
\text{wall crossing and contractions}
\longrightarrow
\text{toric MMP}
\end{gathered}}
\]
as concrete as possible.  We work over \(\CC\), and most examples begin
with a smooth projective toric variety so that
\(\Cl(X)=\Pic(X)\) is free and the Cox group is itself a torus.
\end{chapteroverview}

\begin{threedictionaries}
The same data will be read in three languages throughout the chapter.

\medskip

\begin{center}
\small
\begin{tabularx}{\linewidth}{@{}
  >{\raggedright\arraybackslash}X
  >{\raggedright\arraybackslash}X
  >{\raggedright\arraybackslash}X@{}}
\toprule
\textbf{Cox coordinates}
&
\textbf{GIT}
&
\textbf{Birational geometry}
\\
\midrule
class \(\theta\in\Cl(X)\)
& character / linearization
& choice of model and polarization
\\
homogeneous monomial
& semi-invariant section
& witness for semistability
\\
region of characters
& constant semistable locus
& one toric birational model
\\
wall
& limiting quotient
& contraction or flip
\\
\bottomrule
\end{tabularx}
\end{center}
\end{threedictionaries}

\begin{surfacelaboratory}
Our running two-dimensional example will be
\[
\mathbb F_1\simeq\Bl_p(\PP^2).
\]
Chapter~6 constructed it from \(\AAA^4\) by a \((\Gm)^2\)-action.  In
this chapter we will \emph{keep that affine space and that action fixed}
and vary only the character used to form the quotient.  One region of
characters gives \(\mathbb F_1\); another gives \(\PP^2\); on the wall
between them we will see the exceptional curve of the blowup contract.

This is the surface model for the entire chapter.  In dimension three
the same kind of wall crossing can also produce small contractions and
flips, which have no surface analogue.
\end{surfacelaboratory}

\section{One Cox quotient, many linearizations}
\label{sec:one-cox-many-linearizations}
\index{geometric invariant theory}
\index{linearization}
\index{semistable point}
\index{weight cone}
\index{variation of GIT quotient}

Recall the Cox construction for a smooth projective toric variety
\(X_\Sigma\).  If the rays are
\(\Sigma(1)=\{\rho_1,\dots,\rho_r\}\), then
\[
S=\CC[x_1,\dots,x_r]
\]
is graded by \(\Cl(X_\Sigma)\), with
\[
\deg(x_i)=[D_{\rho_i}].
\]
The Cox group
\[
G=\Hom\bigl(\Cl(X_\Sigma),\Gm\bigr)
\]
acts diagonally on \(\AAA^r\).  Chapter~6 then formed
\[
X_\Sigma
\simeq
\bigl(\AAA^r\setminus Z(\Sigma)\bigr)/G.
\]

There are three ingredients in this display:
\[
\boxed{
\text{the affine space}
\quad+\quad
\text{the action of }G
\quad+\quad
\text{the open set we quotient}.}
\]
The first two are encoded by the Cox grading.  GIT will tell us how to
produce the third from one additional choice: a character of \(G\).

\begin{guidingquestion}
What changes if we keep \(\AAA^r\) and the action of \(G\) fixed, but
change the character used to form the quotient?
\end{guidingquestion}

\subsection{Why the affine quotient is not enough}

If a reductive group \(G\) acts on an affine variety \(V\), the affine
categorical quotient is
\[
V\mathbin{//}G=\Spec\bigl(\CC[V]^G\bigr).
\]
For Cox constructions this often forgets far too much.  The simplest
example is already projective space: scalar multiplication by \(\Gm\)
on \(\AAA^{n+1}\) has only constant polynomial invariants, and hence
\[
\AAA^{n+1}\mathbin{//}\Gm=\{\mathrm{pt}\}.
\]
But deleting the origin and remembering how functions \emph{transform},
rather than insisting that they be invariant, recovers \(\PP^n\).

This is the elementary idea behind the GIT linearization.  Let
\[
\widehat G=\Hom(G,\Gm)
\]
be the character group.  A character \(\chi\in\widehat G\) lifts the
action on \(\AAA^r\) to the trivial line bundle by
\[
g\cdot(z,t)=\bigl(g\cdot z,\chi(g)t\bigr).
\]
Although the underlying line bundle is still trivial, its \(G\)-action
is not.  The choice of \(\chi\) is called a \emph{linearization}.

For the Cox group there is a canonical identification
\[
\boxed{
\widehat G\simeq\Cl(X_\Sigma).}
\]
Thus a divisor class can be read as a character.  If
\(\theta\in\Cl(X_\Sigma)\) and \(\chi_\theta\) is the corresponding
character, then a Cox polynomial \(F\) of degree \(\theta\) satisfies
\[
F(g\cdot z)=\chi_\theta(g)F(z).
\]
Such an \(F\) is a \emph{semi-invariant} of weight \(\chi_\theta\).
This is precisely the relation between the Cox grading and GIT used in
Cox--Little--Schenck, Example~14.1.4.

\begin{agcheatsheet}
For the diagonal torus actions used in this chapter, it is enough to
remember the following translation.
\[
\begin{array}{ccl}
F\in S_{d\theta}
&\Longleftrightarrow&
F(gz)=\chi_\theta(g)^dF(z),
\\[1mm]
F(z)\neq0
&\Longrightarrow&
z\text{ is allowed by the linearization }\theta.
\end{array}
\]
A point is semistable when \emph{some} positive multiple \(d\theta\)
has such a homogeneous polynomial not vanishing there.  So in the Cox
setting, semistability is a question about which homogeneous monomials
can be nonzero at a point.
\end{agcheatsheet}

\subsection{Semistable points and the GIT quotient}

We now make that rule precise.

\begin{definition}[Semistability for a Cox linearization]
\label{def:cox-git-semistable}
Fix \(\theta\in\Cl(X_\Sigma)\).  A point \(z\in\AAA^r\) is
\emph{semistable with respect to \(\theta\)} if there exist an integer
\(d>0\) and a polynomial
\[
F\in S_{d\theta}
\]
such that \(F(z)\neq0\).  The semistable locus is denoted
\((\AAA^r)^{ss}_\theta\).
\end{definition}

This is the toric form of Cox--Little--Schenck,
Definition~14.1.6.  Their definition is phrased in terms of invariant
sections of the linearized bundle; the Cox grading identifies those
sections with the homogeneous pieces \(S_{d\theta}\).

The quotient is built from all of these pieces at once.

\begin{definition}[The toric GIT quotient]
\label{def:toric-git-quotient}
Set
\[
R_\theta
=
\bigoplus_{d\geq0}S_{d\theta}.
\]
The GIT quotient of \(\AAA^r\) by \(G\) with linearization \(\theta\)
is
\[
\boxed{
\AAA^r\mathbin{//}_\theta G
:=
\Proj(R_\theta).}
\]
Equivalently,
\[
\AAA^r\mathbin{//}_\theta G
\simeq
(\AAA^r)^{ss}_\theta\mathbin{//}G.
\]
\end{definition}

\noindent
\textit{Reference.}
This is Cox--Little--Schenck, Definition~14.1.11 and
Proposition~14.1.12.

The point of the construction is that \(G\) and its action have not
changed.  Only \(\theta\) has changed, but \(\theta\) changes the graded
ring \(R_\theta\) and the open set \((\AAA^r)^{ss}_\theta\).  The
quotient can therefore change as well.

For a diagonal torus action there is a remarkably efficient test for
semistability.  Write the weight of the coordinate \(x_i\) as
\[
\beta_i=\deg(x_i)\in\widehat G,
\]
and tensor the character lattice with \(\RR\).

\begin{proposition}[The weight-cone test]
\label{prop:weight-cone-semistability}
Let \(G\simeq(\Gm)^k\) act diagonally on \(\AAA^r\) with weights
\(\beta_1,\dots,\beta_r\), and let \(\theta\in\widehat G\).  For
\(z=(z_1,\dots,z_r)\), set
\[
I(z)=\{i:z_i\neq0\}.
\]
Then
\[
\boxed{
z\in(\AAA^r)^{ss}_\theta
\quad\Longleftrightarrow\quad
\theta\in
\Cone{\beta_i:i\in I(z)}
\subseteq\widehat G_\RR.}
\]
\end{proposition}

\begin{proof}
Suppose first that \(z\) is semistable.  Some
\(F\in S_{d\theta}\) does not vanish at \(z\).  At least one monomial
\[
x^a=x_1^{a_1}\cdots x_r^{a_r}
\]
appearing in \(F\) is nonzero at \(z\).  Hence \(a_i=0\) whenever
\(z_i=0\), and homogeneity gives
\[
d\theta
=
\sum_{i\in I(z)}a_i\beta_i.
\]
After dividing by \(d\), the character \(\theta\) lies in the stated
cone.

Conversely, if \(\theta\) lies in this rational polyhedral cone, then
after clearing denominators there are \(d>0\) and integers
\(a_i\geq0\), supported on \(I(z)\), such that
\[
d\theta=\sum_i a_i\beta_i.
\]
The monomial \(x^a\in S_{d\theta}\) does not vanish at \(z\), so \(z\)
is semistable.
\end{proof}

This test is the computational engine for variation of toric GIT.  The
semistable locus can change only when \(\theta\) moves across one of the
cones spanned by subsets of the weights.

\begin{example}[Scalar multiplication revisited]
Let \(G=\Gm\) act on \(\AAA^{n+1}\) with every coordinate of weight
\(1\), and let \(\theta=d\in\widehat G\simeq\ZZ\).  The weight-cone
test gives
\[
\begin{array}{c|c|c}
d& (\AAA^{n+1})^{ss}_d&\AAA^{n+1}\mathbin{//}_d\Gm\\
\hline
d>0&\AAA^{n+1}\setminus\{0\}&\PP^n\\
d=0&\AAA^{n+1}&\{\mathrm{pt}\}\\
d<0&\varnothing&\varnothing.
\end{array}
\]
Thus even the usual construction of projective space already exhibits
the basic VGIT phenomenon: the quotient depends on the character, not
only on the action.
\end{example}

\subsection{An ample class recovers the Cox quotient}

The original Cox construction has not disappeared inside this larger
family.  It is recovered by choosing a positive character.

\begin{proposition}[Cox quotient as a GIT quotient]
\label{prop:ample-character-recovers-cox}
Let \(X_\Sigma\) be a smooth projective toric variety and let
\(\theta\in\Cl(X_\Sigma)\) be ample.  Then
\[
\boxed{
(\AAA^{\Sigma(1)})^{ss}_\theta
=
\AAA^{\Sigma(1)}\setminus Z(\Sigma),}
\]
and
\[
\AAA^{\Sigma(1)}\mathbin{//}_\theta G
\simeq X_\Sigma.
\]
In this smooth case the semistable points are stable and the quotient
is geometric.
\end{proposition}

\noindent
\textit{Reference.}
The equality of the semistable locus with the Cox open set is
Cox--Little--Schenck, Proposition~14.1.9.  In their more general
projective setting, the stable locus is obtained from the subfan of
simplicial cones; for a smooth fan this is all of \(\Sigma\).

This proposition changes the way we should think about Chapter~6.
There we started from \(\Sigma\), built \(Z(\Sigma)\), and then took a
quotient.  Now we can start from the diagonal Cox action and an ample
class \(\theta\); semistability reconstructs the deleted set for us.
Better still, \(\theta\) is free to move.

\subsection{Surface laboratory: the same action gives
\texorpdfstring{\(\mathbb F_1\)}{F1} and
\texorpdfstring{\(\PP^2\)}{P2}}

Return to the Cox action for
\[
X=\mathbb F_1\simeq\Bl_p(\PP^2)
\]
from Chapter~6.  With Cox coordinates
\[
(u_0,v_0,u_1,v_1)
\]
and the basis \((F,S)\) of \(\Cl(X)\simeq\ZZ^2\), the weights are
\[
\deg(u_0)=\deg(u_1)=F,
\qquad
\deg(v_0)=S,
\qquad
\deg(v_1)=F+S.
\]
Equivalently,
\[
(\lambda,\mu)\cdot(u_0,v_0,u_1,v_1)
=
(\lambda u_0,\mu v_0,\lambda u_1,\lambda\mu v_1).
\]

Now keep this action fixed and vary
\[
\theta=aF+bS\in\Cl(X)_\RR.
\]
The three distinct weight rays divide the positive quadrant into two
regions.

\begin{figure}[ht]
\centering
\begin{tikzpicture}[scale=1.18]
  \fill[DictionaryBlue!5] (0,0)--(4.3,0)--(0,4.3)--cycle;
  \fill[PredictionTeal!12] (0,0)--(4.3,0)--(4.3,4.3)--cycle;
  \fill[ExampleGold!12] (0,0)--(4.3,4.3)--(0,4.3)--cycle;

  \draw[->,thick] (-0.2,0)--(4.75,0) node[right] {\(F\)};
  \draw[->,thick] (0,-0.2)--(0,4.75) node[above] {\(S\)};
  \draw[ray,very thick,DictionaryBlue] (0,0)--(4.25,4.25)
    node[above right] {\(F+S\)};

  \node[PredictionTeal!80!black,align=center] at (2.8,1.15)
    {\(a>b>0\)\\\(\mathbb F_1\)};
  \node[ExampleGold!90!black,align=center] at (1.15,2.8)
    {\(b>a>0\)\\\(\PP^2\)};

  \fill (3.7,0) circle (1.4pt);
  \node[below] at (3.7,0) {weight \(F\)};
  \fill (0,3.7) circle (1.4pt);
  \node[left] at (0,3.7) {weight \(S\)};
\end{tikzpicture}
\caption{The character space for the \(\mathbb F_1\) Cox action.  The
diagonal ray already predicts the wall between two generic GIT
quotients.  Section~10.2 will turn this picture into the secondary fan.}
\label{fig:f1-first-vgit-chambers}
\end{figure}

First suppose
\[
a>b>0.
\]
The character lies between the rays \(F\) and \(F+S\).  By the
weight-cone test, a semistable point must have at least one of
\(u_0,u_1\) nonzero and at least one of \(v_0,v_1\) nonzero.  Conversely
those two conditions are sufficient.  Hence
\[
(\AAA^4)^{ss}_\theta
=
\AAA^4\setminus
\bigl(V(u_0,u_1)\cup V(v_0,v_1)\bigr).
\]
This is exactly the Cox open set of \(\mathbb F_1\), so
\[
\boxed{
\AAA^4\mathbin{//}_\theta(\Gm)^2\simeq\mathbb F_1
\qquad(a>b>0).}
\]
In the blowup notation of Chapter~6, \(F=H-E\) and \(S=E\), so
\[
aF+bS=aH-(a-b)E.
\]
The inequalities \(a>b>0\) are exactly the familiar ample inequalities
for the blowup.

Now move \(\theta\) across the diagonal and suppose
\[
b>a>0.
\]
The weight-cone test gives a different open set:
\[
(\AAA^4)^{ss}_\theta
=
\{v_0\neq0\}
\cap
\{(u_0,u_1,v_1)\neq(0,0,0)\}.
\]
Indeed, the weight \(S\) is now indispensable: the cone generated by
\(F\) and \(F+S\) lies on or below the diagonal.  Once \(v_0\neq0\),
semistability requires one additional nonzero coordinate of positive
\(F\)-weight.

Use the \(\mu\)-action to set \(v_0=1\).  The residual \(\lambda\)-action
on the remaining coordinates is
\[
\lambda\cdot(u_0,u_1,v_1)
=(\lambda u_0,\lambda u_1,\lambda v_1).
\]
Therefore
\[
\boxed{
\AAA^4\mathbin{//}_\theta(\Gm)^2\simeq\PP^2
\qquad(b>a>0).}
\]

Nothing about \(\AAA^4\) changed.  Nothing about the
\((\Gm)^2\)-action changed.  The character crossed one ray, the
semistable locus changed, and the quotient changed from the blowup to
the projective plane.

\subsection{What lives on the wall?}

The separating ray is
\[
\theta=F+S=H.
\]
Its degree-one semi-invariants are
\[
u_0v_0,
\qquad
u_1v_0,
\qquad
v_1.
\]
In fact every monomial of degree \(d(F+S)\) is a product of these
three, so
\[
R_{F+S}
\simeq
\CC[u_0v_0,u_1v_0,v_1]
\]
with all three generators in degree one.  Consequently the wall
quotient is again
\[
\AAA^4\mathbin{//}_{F+S}(\Gm)^2\simeq\PP^2.
\]

From the \(\mathbb F_1\)-side, the quotient map is visibly given by
\[
\mathbb F_1\longrightarrow\PP^2,
\qquad
[u_0,v_0,u_1,v_1]
\longmapsto
[u_0v_0:u_1v_0:v_1].
\]
On the divisor \(v_0=0\), every point is sent to
\([0:0:1]\).  Under the Chapter~6 identification \(S=E\), this is the
exceptional curve of \(\Bl_p(\PP^2)\).  We have therefore rediscovered
the blowdown
\[
\boxed{
\mathbb F_1\simeq\Bl_p(\PP^2)
\longrightarrow
\PP^2}
\]
by moving a character to a wall.

\begin{pauseandthink}
The equation \(a=b\) did not come from the fan of \(\mathbb F_1\) by
inspection.  It came from the arrangement of the \emph{Cox weights} in
\(\Cl(X)_\RR\).  Why should all characters with \(a>b>0\) give the same
semistable locus?  Why should all characters with \(b>a>0\) do the
same?  What finite collection of cones in character space organizes
this behavior in general?
\end{pauseandthink}


The answer is the secondary fan.  Its chambers are precisely the
regions on which the toric GIT quotient is constant in the appropriate
sense.  The example above has already drawn the relevant part of that
fan; the next section will construct it systematically and explain its
relation with regular subdivisions and the secondary polytope.

\section{The secondary fan and the secondary polytope}
\label{sec:secondary-fan-polytope}
\index{secondary fan}
\index{GKZ fan|see{secondary fan}}
\index{Gale duality}
\index{regular triangulation}
\index{secondary polytope}

Section~\ref{sec:one-cox-many-linearizations} found a wall by hand.  The
weights for the Cox action on \(\mathbb F_1\) were
\[
F,\qquad S,\qquad F,\qquad F+S,
\]
and the ray through \(F+S\) separated two regions with different
quotients.  Nothing in that computation, however, explained why there
should be only finitely many such regions or why they should fit
together as a fan.

The answer is remarkably rigid.  For a diagonal torus action, the
weights themselves cut character space into a rational polyhedral fan.
This is the \emph{secondary fan}, also called the \emph{GKZ
decomposition}.  Its chambers record the generic GIT quotients, while
its faces record the quotients that appear when the character becomes
special.

\begin{guidingquestion}
How can we read \emph{all} toric GIT quotients of a fixed diagonal
action from one finite combinatorial object?
\end{guidingquestion}

\subsection{From weight cones to chambers}

Let
\[
G\simeq(\Gm)^k
\]
act diagonally on \(\AAA^r\) with weights
\[
\beta_1,\dots,\beta_r\in\widehat G,
\]
and set
\[
C_\beta
=
\Cone{\beta_1,\dots,\beta_r}
\subseteq\widehat G_\RR.
\]
The weight-cone test of
Proposition~\ref{prop:weight-cone-semistability} immediately shows that
characters outside \(C_\beta\) have empty semistable locus.  Thus
\(C_\beta\) is the region in character space where nonempty GIT
quotients can occur.

For a point \(z=(z_1,\dots,z_r)\), recall
\[
I(z)=\{i:z_i\neq0\}.
\]
It is useful to package exactly the weights visible at \(z\).

\begin{definition}[Weight cone of a point]
The \emph{weight cone} of \(z\in\AAA^r\) is
\[
\omega(z)
=
\Cone{\beta_i:i\in I(z)}
\subseteq\widehat G_\RR.
\]
Thus
\[
z\in(\AAA^r)^{ss}_\theta
\quad\Longleftrightarrow\quad
\theta\in\omega(z).
\]
\end{definition}

There are only finitely many possible cones \(\omega(z)\), since each is
spanned by a subset of the finite set \(\{\beta_1,\dots,\beta_r\}\).
For \(\theta\in C_\beta\), define
\[
\Gamma(\theta)
=
\bigcap_{\theta\in\omega(z)}\omega(z).
\]
In words: intersect every weight cone in which \(\theta\) lies.  The
result is the smallest region that the semistability test can detect
around \(\theta\).

\begin{theorem}[The secondary fan]
\label{thm:secondary-fan}
The cones \(\Gamma(\theta)\), together with their faces, form a finite
rational polyhedral fan
\[
\Sigma_{\mathrm{GKZ}}
\quad\text{in}\quad
\widehat G_\RR
\]
whose support is \(C_\beta\).  Moreover, if \(\theta\) and \(\theta'\)
lie in the relative interior of the same cone, then
\[
(\AAA^r)^{ss}_\theta
=
(\AAA^r)^{ss}_{\theta'}
\]
and hence
\[
\AAA^r\mathbin{//}_\theta G
\simeq
\AAA^r\mathbin{//}_{\theta'}G.
\]
The maximal cones of \(\Sigma_{\mathrm{GKZ}}\) are called
\emph{chambers}; their codimension-one faces are called \emph{walls}.
\end{theorem}

\noindent
\textit{Reference.}
This is the diagonal-torus form of Cox--Little--Schenck,
Theorem~14.4.7.  Their GKZ cones keep track simultaneously of the fan
of the quotient and of coordinates that never vanish; the description
above is the same chamber structure viewed directly through
semistability.

The theorem is substantial, but its basic mechanism is already visible
from the weight-cone test.  For every point \(z\), the statement
``\(z\) is semistable'' is the polyhedral condition
\[
\theta\in\omega(z).
\]
Since only finitely many \(\omega(z)\) occur, only finitely many
semistability patterns can occur.  Passing from this finite collection
of cones to their common chamber decomposition gives
\(\Sigma_{\mathrm{GKZ}}\).  The content of the theorem is that these
regions fit together with the face-to-face behavior required of a fan.

\begin{agcheatsheet}
For computations, one can think of the secondary fan as a
\emph{master map of linearizations}:
\[
\boxed{
\begin{array}{ccl}
\theta\notin C_\beta
&\Longrightarrow&
\text{empty semistable locus},
\\[1mm]
\theta\text{ in a chamber}
&\Longrightarrow&
\text{one generic toric quotient},
\\[1mm]
\theta\text{ on a wall}
&\Longrightarrow&
\text{a limiting quotient},
\\[1mm]
\theta\text{ on a lower-dimensional face}
&\Longrightarrow&
\text{a still more special quotient}.
\end{array}}
\]
The important point is not merely that the quotient is constant inside
a chamber.  The \emph{way the chambers meet} records which models can
be connected by varying the linearization.
\end{agcheatsheet}

\subsection{The \texorpdfstring{\(\mathbb F_1\)}{F1} picture really is a fan}

Return to the action from the previous section.  The distinct weight
rays are
\[
\RR_{\geq0}F,
\qquad
\RR_{\geq0}(F+S),
\qquad
\RR_{\geq0}S,
\]
and
\[
C_\beta=\Cone{F,S}.
\]
There are two two-dimensional GKZ chambers:
\[
\Gamma_{\mathbb F_1}
=
\Cone{F,F+S},
\qquad
\Gamma_{\PP^2}
=
\Cone{F+S,S}.
\]
Indeed, if \(\theta=aF+bS\), then
\[
\begin{array}{c|c|c}
\text{location of }\theta
&\text{semistable locus}
&\text{quotient}
\\
\hline
a>b>0
&\AAA^4\setminus
\bigl(V(u_0,u_1)\cup V(v_0,v_1)\bigr)
&\mathbb F_1
\\[1mm]
a=b>0
&\text{wall semistable locus}
&\PP^2
\\[1mm]
b>a>0
&\{v_0\neq0\}\cap
\{(u_0,u_1,v_1)\neq(0,0,0)\}
&\PP^2.
\end{array}
\]

Thus Figure~\ref{fig:f1-first-vgit-chambers} was not just a helpful
picture: it was the two-dimensional part of
\(\Sigma_{\mathrm{GKZ}}\).  The common face
\[
\RR_{\geq0}(F+S)
\]
is the wall across which the exceptional curve of \(\mathbb F_1\) is
contracted.

There is a small subtlety worth remembering.  The chamber
\(\Gamma_{\mathbb F_1}\) is literally
\[
\Nef(\mathbb F_1)
=
\Cone{H-E,H}
=
\Cone{F,F+S}.
\]
On the other side, however, the quotient is \(\PP^2\), whose Picard
group has rank one even though \(\Gamma_{\PP^2}\) is two-dimensional.
The reason is visible in the semistable locus: \(v_0\) is forced to be
nonzero, so one Cox coordinate has disappeared from the quotient fan.
The secondary fan remembers this extra GIT information.  We will return
to the relation with nef cones in Section~10.4.

\subsection{Where triangulations enter: Gale duality}

So far the secondary fan lives in character space.  Its second life is
as a space of \emph{regular triangulations}.  The bridge is Gale
duality, the same linear algebra that relates Cox weights to ray
generators.

For the cleanest geometric picture, suppose that
\[
\nu_1,\dots,\nu_r\in N
\]
are distinct lattice points lying in an affine hyperplane
\(H\subset N_\RR\) with \(0\notin H\), and that
\[
Q_\nu=\Conv(\nu_1,\dots,\nu_r)
\]
has full dimension in \(H\).  The map
\[
M\longrightarrow\ZZ^r,
\qquad
m\longmapsto
\bigl(\langle m,\nu_1\rangle,\dots,
      \langle m,\nu_r\rangle\bigr)
\]
has a free cokernel after making the usual lattice hypotheses.  Write
that cokernel as \(\widehat G\).  After tensoring with \(\RR\), we have
\[
0\longrightarrow M_\RR
\longrightarrow\RR^r
\longrightarrow\widehat G_\RR
\longrightarrow0.
\]
The images \(\beta_i\) of the standard basis vectors are precisely the
Gale-dual weights.

Now assign a height \(\omega_i\geq0\) to every \(\nu_i\) and form the
lifted cone
\[
C_{\nu,\omega}
=
\Cone{(\nu_1,\omega_1),\dots,(\nu_r,\omega_r)}
\subseteq N_\RR\times\RR.
\]
Project the lower faces back to \(N_\RR\).  Their cones give a
subdivision \(\Sigma_\omega\) of \(C_\nu\).  Intersecting with the
hyperplane \(H\) gives a subdivision of \(Q_\nu\).

\begin{definition}[Regular triangulation]
A triangulation \(\mathcal T\) of \(Q_\nu\), using vertices among the
\(\nu_i\), is \emph{regular} (or \emph{coherent}) if it is obtained by
this lifting-and-lower-hull construction for some choice of heights
\(\omega=(\omega_1,\dots,\omega_r)\).
\end{definition}

This is Cox--Little--Schenck, Definition~15.2.8.  The construction
explains why the height vector naturally lives in the same space as a
GIT character.  If we change \(\omega\) by
\[
\omega_i
\longmapsto
\omega_i+\langle m,\nu_i\rangle,
\]
we have merely tilted the lifting by a linear function.  The regular
subdivision does not change.  Hence the meaningful parameter is the
class
\[
[\omega]\in
\RR^r/M_\RR
\simeq
\widehat G_\RR.
\]
That quotient is exactly the character space in which the secondary
fan lives.

\begin{theorem}[Chambers and regular triangulations]
\label{thm:gkz-chambers-triangulations}
In the affine-hyperplane situation above, the chambers of
\(\Sigma_{\mathrm{GKZ}}\) are in bijection with the regular
triangulations of \(\nu\).  A character in a chamber determines the
corresponding triangulation by the lower-hull construction.
\end{theorem}

\noindent
\textit{Reference.}
This is Cox--Little--Schenck, Proposition~15.2.9.  We will use the
theorem as a dictionary rather than reproduce its full proof: Gale
duality identifies height vectors modulo linear functions with GIT
characters, while strict convexity of the resulting piecewise linear
function is exactly the condition that the projected lower hull be a
triangulation.

Thus, in this setting, the same chamber can be read in three ways:
\[
\boxed{
\begin{gathered}
\text{GIT chamber in }\widehat G_\RR
\quad\longleftrightarrow\quad
\text{regular triangulation of }Q_\nu
\\[-1mm]
\longleftrightarrow
\text{simplicial fan of the corresponding toric quotient}.
\end{gathered}}
\]

\subsection{The secondary polytope}

There is now a beautiful dual way to package all the chambers at once.
For a triangulation \(\mathcal T\), define its \emph{GKZ vector}
\[
\phi_{\mathcal T}
=
(\phi_1,\dots,\phi_r)\in\RR^r,
\]
where
\[
\phi_i
=
\sum_{\substack{\Delta\in\mathcal T\\ \nu_i\in\Delta}}
\vol(\Delta).
\]
Here \(\Delta\) runs over the maximal simplices and
\(\vol\) denotes
normalized lattice volume in \(H\).  The classical GKZ secondary
polytope is the convex hull of these vectors; equivalently, after
passing to its natural affine span, it is a polytope dual to the
secondary fan.

\begin{theorem}[The secondary polytope]
\label{thm:secondary-polytope}
Under the affine-hyperplane hypotheses above, there is a full
dimensional lattice polytope
\[
P_{\mathrm{GKZ}}\subseteq\widehat G_\RR^*
\]
whose normal fan is \(\Sigma_{\mathrm{GKZ}}\).  Its vertices are in
bijection with regular triangulations of \(\nu\).
\end{theorem}

\noindent
\textit{Reference.}
See Cox--Little--Schenck, Proposition~15.2.11, for the existence and
normal-fan statement.  The vertex--triangulation correspondence is
Corollary~15.2.12.  The GKZ-vector construction above gives the
classical distinguished realization of the secondary polytope.  Our
normal-fan convention may be replaced by the opposite convention if
one uses upper rather than lower hulls; none of the chamber
combinatorics changes.

The word \emph{secondary} is now rather apt.  The original polytope
\(Q_\nu\) lives in the lattice space and admits many triangulations.
The secondary polytope lives in the dual parameter space and organizes
\emph{the triangulations themselves}.

\begin{center}
\small
\begin{tabular}{c|c}
\textbf{Secondary polytope}
&\textbf{Secondary fan / VGIT}
\\
\hline
vertex
&chamber / regular triangulation
\\
edge
&wall between adjacent chambers
\\
higher-dimensional face
&coarser regular subdivision
\\
normal cone to a face
&characters producing that subdivision
\end{tabular}
\end{center}

The reversal of dimensions is exactly what one expects from a normal
fan: vertices of a polytope have full-dimensional normal cones, while
edges correspond to codimension-one cones.  Thus a path from one
vertex of \(P_{\mathrm{GKZ}}\) to a neighboring vertex is dual to
crossing a wall in the VGIT fan.

\subsection{A square: two triangulations and one wall}

The smallest nontrivial secondary polytope can be computed completely
by hand.  Take the four lattice points
\[
\nu_1=(0,0,1),\quad
\nu_2=(1,0,1),\quad
\nu_3=(1,1,1),\quad
\nu_4=(0,1,1)
\]
in the affine plane \(z=1\).  Their convex hull is a square.  There are
exactly two triangulations:
\[
\mathcal T_{13}
=
\{123,134\},
\qquad
\mathcal T_{24}
=
\{124,234\},
\]
according to which diagonal is chosen.

The unique affine dependence is
\[
\nu_1-\nu_2+\nu_3-\nu_4=0.
\]
Consequently the one-dimensional quotient of height space is measured
by
\[
\delta(\omega)
=
\omega_1+\omega_3-\omega_2-\omega_4.
\]
If \(\delta(\omega)<0\), the points \(\nu_1,\nu_3\) lie relatively
lower and the lower hull produces \(\mathcal T_{13}\).  If
\(\delta(\omega)>0\), it produces \(\mathcal T_{24}\).  When
\(\delta(\omega)=0\), the four lifted points are coplanar and the
projected lower face is the entire square: the diagonal disappears.

\begin{figure}[ht]
\centering
\begin{tikzpicture}[scale=1.02]
  \coordinate (A1) at (0,0);
  \coordinate (A2) at (2.2,0);
  \coordinate (A3) at (2.2,2.2);
  \coordinate (A4) at (0,2.2);
  \fill[DictionaryBlue!7] (A1)--(A2)--(A3)--cycle;
  \fill[PredictionTeal!9] (A1)--(A3)--(A4)--cycle;
  \draw[thick] (A1)--(A2)--(A3)--(A4)--cycle;
  \draw[very thick,DictionaryBlue] (A1)--(A3);
  \node[below left] at (A1) {\(1\)};
  \node[below right] at (A2) {\(2\)};
  \node[above right] at (A3) {\(3\)};
  \node[above left] at (A4) {\(4\)};
  \node[align=center] at (1.1,-0.65)
    {\(\mathcal T_{13}\)\\\(\delta<0\)};

  \draw[very thick,ExampleGold!90!black] (3.45,1.1)--(5.95,1.1);
  \fill[ExampleGold!90!black] (3.45,1.1) circle (2.2pt);
  \fill[ExampleGold!90!black] (5.95,1.1) circle (2.2pt);
  \node[above,align=center] at (4.70,1.1)
    {secondary polytope\\\(P_{\mathrm{GKZ}}\)};
  \node[below] at (3.45,1.0) {\(\phi_{13}\)};
  \node[below] at (5.95,1.0) {\(\phi_{24}\)};

  \coordinate (B1) at (7.2,0);
  \coordinate (B2) at (9.4,0);
  \coordinate (B3) at (9.4,2.2);
  \coordinate (B4) at (7.2,2.2);
  \fill[PredictionTeal!9] (B1)--(B2)--(B4)--cycle;
  \fill[DictionaryBlue!7] (B2)--(B3)--(B4)--cycle;
  \draw[thick] (B1)--(B2)--(B3)--(B4)--cycle;
  \draw[very thick,DictionaryBlue] (B2)--(B4);
  \node[below left] at (B1) {\(1\)};
  \node[below right] at (B2) {\(2\)};
  \node[above right] at (B3) {\(3\)};
  \node[above left] at (B4) {\(4\)};
  \node[align=center] at (8.3,-0.65)
    {\(\mathcal T_{24}\)\\\(\delta>0\)};
\end{tikzpicture}
\caption{The square has two regular triangulations, hence its secondary
polytope is a segment.  The wall \(\delta=0\) is the coarser
subdivision consisting of the unsplit square.}
\label{fig:square-secondary-polytope}
\end{figure}

Because each triangle has normalized area one, the two GKZ vectors are
\[
\phi_{13}=(2,1,2,1),
\qquad
\phi_{24}=(1,2,1,2).
\]
Thus
\[
P_{\mathrm{GKZ}}
=
\Conv(\phi_{13},\phi_{24})
\]
is a line segment.  Its normal fan is one-dimensional: it consists of
two rays separated by the origin.  Those two chambers are exactly
\[
\delta<0
\qquad\text{and}\qquad
\delta>0,
\]
with the wall \(\delta=0\) between them.  If one defines normal cones
using maximization rather than the lower-hull/minimization convention,
the two ray labels are exchanged; the fan itself is unchanged.

This tiny example contains the entire VGIT pattern.  A generic
parameter chooses one triangulation.  Moving to the wall erases a
diagonal and produces a coarser subdivision.  Crossing the wall inserts
the other diagonal.  In higher dimension this elementary replacement
is called a \emph{bistellar flip}; it is the combinatorial shadow of the
birational transformations we will study next.

\begin{pauseandthink}
The secondary fan and the secondary polytope contain the same
incidence data, but with dimensions reversed.  A chamber of the fan is
dual to a vertex of the polytope, and a wall between two chambers is
dual to an edge between two vertices.  Thus there are two equivalent
ways to imagine varying a linearization: move \emph{through} the
secondary fan, or move \emph{along edges} of the secondary polytope.
Which picture makes a wall crossing easier to see?
\end{pauseandthink}

We are now ready to put geometry back into the picture.  A wall of the
secondary fan tells us that two GIT chambers are adjacent; a flip of
regular triangulations tells us how their combinatorics differ.  The
next question is what the induced map between the corresponding toric
varieties actually does.  In Section~10.3 we will see that, depending
on the wall, it can contract a divisor, produce a small contraction and
a flip, or lead to a fiber-type morphism.

\section{Crossing a wall: contractions, flips, and fibrations}
\label{sec:crossing-wall}
\index{wall crossing}
\index{contraction!divisorial}
\index{contraction!small}
\index{contraction!fiber type}
\index{flip!toric}
\index{flop!toric}

The secondary fan tells us \emph{when} the quotient changes.  We now
ask what the change means geometrically.  Let \(C_-\) and \(C_+\) be
adjacent chambers, choose characters
\[
\theta_-\in\relint(C_-),
\qquad
\theta_+\in\relint(C_+),
\]
and let \(\theta_0\) lie in the relative interior of their common wall.
Write
\[
X_-
=
\AAA^r\mathbin{//}_{\theta_-}G,
\qquad
X_0
=
\AAA^r\mathbin{//}_{\theta_0}G,
\qquad
X_+
=
\AAA^r\mathbin{//}_{\theta_+}G.
\]
Under the usual semiprojective hypotheses, variation of GIT produces
projective toric morphisms
\[
\begin{tikzcd}[column sep=huge]
X_-
  \arrow[dr,"\pi_-"']
  \arrow[rr,dashed,"\psi"]
&&
X_+
  \arrow[dl,"\pi_+"]
\\
&X_0.&
\end{tikzcd}
\]
The wall quotient \(X_0\) forgets just enough information that the two
nearby models can map to it.  The dashed arrow is the induced
birational map when both morphisms are birational.

\begin{guidingquestion}
Which subvarieties become invisible when a character reaches a wall,
and how does the answer determine the birational transformation on the
other side?
\end{guidingquestion}

\subsection{The three geometric possibilities}

The geometry is controlled by the exceptional locus.  Recall that the
exceptional locus of a birational morphism is the set on which it is
not an isomorphism.

\begin{definition}[Three types of contraction]
Let \(\pi:X\to Y\) be a projective toric contraction with connected
fibers.
\begin{enumerate}
\item It is of \emph{fiber type} if
      \(\dim Y<\dim X\).
\item It is \emph{divisorial} if it is birational and its exceptional
      locus contains a divisor.
\item It is \emph{small} if it is birational and its exceptional locus
      has codimension at least two.
\end{enumerate}
If two small contractions
\[
X_-\longrightarrow X_0\longleftarrow X_+
\]
arise from the two sides of a wall, the birational map
\(X_-\dashrightarrow X_+\) is an \emph{elementary flip}.  It is a
\emph{flop} when the canonical divisor has degree zero on the curves
contracted on either side.
\end{definition}

The adjective \emph{small} refers to the exceptional locus, not to the
size of the fibers.  A small contraction may collapse curves, but it
collapses no divisor.  Consequently it does not delete a ray from the
fan: the fans on the two sides have the same rays and differ only in
how those rays are assembled into cones.

\begin{proposition}[Wall-crossing trichotomy]
\label{prop:wall-crossing-trichotomy}
Let \(X_\Sigma\) be a simplicial semiprojective toric variety, and let a
codimension-one face of its nef chamber determine the wall morphism
\[
\pi:X_\Sigma\longrightarrow X_{\Sigma_0}.
\]
Exactly one of the following occurs:
\[
\begin{array}{c|c|c}
\textbf{type}
&\textbf{geometry of }\pi
&\textbf{fan behavior}
\\
\hline
\text{fiber type}
&\dim X_{\Sigma_0}<\dim X_\Sigma
&\Sigma_0\text{ is degenerate}
\\
\text{divisorial}
&\pi\text{ contracts an invariant divisor}
&\text{one ray is removed}
\\
\text{small}
&\codim\Exc(\pi)\geq2
&\Sigma(1)=\Sigma_0(1).
\end{array}
\]
In the small case, a chamber on the other side supplies a second
simplicial refinement \(\Sigma_+\) of \(\Sigma_0\) and hence an
elementary flip
\[
X_\Sigma\dashrightarrow X_{\Sigma_+}.
\]
\end{proposition}

\noindent
\textit{Reference.}
The classification by exceptional locus is
Cox--Little--Schenck, Proposition~15.4.5.  The wall diagram and the
description of the exceptional loci for a flipping wall are given in
Theorem~15.3.13; see also Lemma~15.3.7.  We postpone the interpretation
in terms of extremal rays of the Mori cone until Section~10.4.

Here is the practical picture to retain.
\[
\boxed{
\begin{array}{ccl}
\text{a whole direction survives in every fiber}
&\Longrightarrow&
\text{fiber type},
\\[1mm]
\text{one Cox ray disappears}
&\Longrightarrow&
\text{divisorial contraction},
\\[1mm]
\text{the rays stay but the cones change}
&\Longrightarrow&
\text{flip or flop}.
\end{array}}
\]

\subsection{Surface laboratory: the two boundary rays of
\texorpdfstring{\(\Nef(\mathbb F_1)\)}{Nef(F1)}}

The chamber for \(\mathbb F_1\) in Section~10.1 was
\[
\Cone{F,H},
\qquad
H=F+S,
\]
with interior classes
\[
\theta=aF+bS,
\qquad
a>b>0.
\]
Its two boundary rays have completely different meanings.  The ray
through \(H\) gives the divisorial blowdown already seen in
Section~10.1; the ray through \(F\) gives the ruling
\(\mathbb F_1\to\PP^1\).

\begin{example}[The wall \(H\): a divisorial contraction]
At \(\theta=H=F+S\), the degree-one semi-invariants are
\[
u_0v_0,
\qquad
u_1v_0,
\qquad
v_1.
\]
They generate the section ring, so the wall quotient is
\[
\Proj(R_H)
\simeq
\PP^2.
\]
The induced morphism is
\[
\pi_H:\mathbb F_1\longrightarrow\PP^2,
\qquad
[u_0,v_0,u_1,v_1]
\longmapsto
[u_0v_0:u_1v_0:v_1].
\]
On the invariant divisor
\[
E=V(v_0),
\]
the first two coordinates vanish, so all of \(E\simeq\PP^1\) maps to
\([0:0:1]\).  Away from \(E\), the map is an isomorphism onto
\(\PP^2\setminus\{[0:0:1]\}\).  Thus
\[
\boxed{
\pi_H:
\Bl_p(\PP^2)\longrightarrow\PP^2
\text{ contracts the divisor }E.}
\]
Combinatorially, crossing this wall removes the ray introduced by the
star subdivision that created the blowup.
\end{example}

\begin{example}[The wall \(F\): a fiber-type contraction]
Now move to the other boundary ray, \(\theta=F\).  The only monomials of
degree \(dF\) are homogeneous polynomials in \(u_0,u_1\).  Therefore
\[
R_F
=
\bigoplus_{d\geq0}\CC[u_0,u_1]_d,
\qquad
\Proj(R_F)\simeq\PP^1.
\]
The morphism defined by this ring is
\[
\pi_F:\mathbb F_1\longrightarrow\PP^1,
\qquad
[u_0,v_0,u_1,v_1]
\longmapsto
[u_0:u_1].
\]
Fixing the point \([u_0:u_1]\) leaves the pair \((v_0,v_1)\), not both
zero, modulo one scalar action.  Hence every fiber is a copy of
\(\PP^1\).  We recover the Hirzebruch-surface ruling
\[
\boxed{
\PP^1\longrightarrow\mathbb F_1
\xrightarrow{\ \pi_F\ }\PP^1.}
\]
Nothing exceptional is confined to a proper subset of \(\mathbb F_1\):
the dimension of the quotient itself drops from two to one.  This is
the fiber-type case.
\end{example}

The intersection pairing gives a second way to recognize the two
maps.  Since
\[
S^2=-1,
\qquad
F^2=0,
\qquad
S\cdot F=1,
\]
we have
\[
H\cdot S=(F+S)\cdot S=0,
\qquad
F\cdot F=0.
\]
The class \(H\) therefore forgets the exceptional curve \(S=E\), while
the class \(F\) has degree zero along every fiber of the ruling.  The
wall in divisor space and the curves contracted by its linear system
are dual pieces of information.  Section~10.4 will turn this
observation into the nef-cone/Mori-cone dictionary.

\begin{center}
\small
\begin{tabular}{c|c|c|c}
\textbf{wall class}
&\textbf{target}
&\textbf{contracted curves}
&\textbf{type}
\\
\hline
\(H=F+S\)
&\(\PP^2\)
&the single curve \(S=E\)
&divisorial
\\
\(F\)
&\(\PP^1\)
&all fibers of the ruling
&fiber type
\end{tabular}
\end{center}

\subsection{Why a genuine flip first appears in dimension three}

A proper birational morphism of normal surfaces cannot be small: every
nontrivial exceptional component is a curve, and a curve on a surface
has codimension one.  Thus a birational wall move between surfaces
must contract a divisor.  To see a small contraction, and therefore a
flip, we need dimension at least three.

The square from Section~10.2 gives the smallest possible example.  Its
ray generators satisfy the circuit relation
\[
\nu_1-\nu_2+\nu_3-\nu_4=0.
\]
The Gale-dual action is therefore the \(\Gm\)-action on \(\AAA^4\)
with weights
\[
(1,-1,1,-1):
\qquad
\lambda\cdot(x_1,x_2,x_3,x_4)
=
(\lambda x_1,\lambda^{-1}x_2,
 \lambda x_3,\lambda^{-1}x_4).
\]
There are two chambers, \(\theta>0\) and \(\theta<0\), separated by the
wall \(\theta=0\).

\begin{example}[The two chamber quotients]
For \(\theta>0\), the weight-cone test says
\[
(\AAA^4)^{ss}_+
=
\{(x_1,x_3)\neq(0,0)\}.
\]
The positive-weight coordinates determine a point
\([x_1:x_3]\in\PP^1\), while \(x_2,x_4\) have weight \(-1\).  Hence
\[
X_+
=
\AAA^4\mathbin{//}_{+}\Gm
\simeq
\operatorname{Tot}
\bigl(\mathcal O_{\PP^1}(-1)^{\oplus2}\bigr).
\]
Similarly,
\[
(\AAA^4)^{ss}_-
=
\{(x_2,x_4)\neq(0,0)\},
\qquad
X_-
\simeq
\operatorname{Tot}
\bigl(\mathcal O_{\PP^1}(-1)^{\oplus2}\bigr),
\]
where the base is now \([x_2:x_4]\).

Both threefolds are abstractly total spaces of the same vector bundle,
but their maps to the wall quotient will remember different
\(\PP^1\)'s.  Under the sign convention of Section~10.2, \(X_+\)
corresponds to the triangulation \(\mathcal T_{24}\), while \(X_-\)
corresponds to \(\mathcal T_{13}\).
\end{example}

\begin{example}[The wall quotient is the quadric cone]
At \(\theta=0\), all four coordinates are semistable and we take the
affine quotient.  The invariant ring is generated by
\[
a=x_1x_2,
\qquad
b=x_1x_4,
\qquad
c=x_3x_2,
\qquad
d=x_3x_4.
\]
These generators satisfy one relation,
\[
ad-bc=0.
\]
Thus
\[
\boxed{
X_0
=
\AAA^4\mathbin{//}_0\Gm
\simeq
Y:=V(ad-bc)\subseteq\AAA^4_{a,b,c,d}.}
\]
The hypersurface \(Y\) is three-dimensional and singular only at the
origin.  It is the affine cone over the quadric surface
\(\PP^1\times\PP^1\subseteq\PP^3\).
\end{example}

The four invariants define morphisms
\[
\begin{tikzcd}[column sep=huge]
X_-
  \arrow[dr,"\pi_-"']
  \arrow[rr,dashed,"\text{flop}"]
&&
X_+
  \arrow[dl,"\pi_+"]
\\
&Y=V(ad-bc).&
\end{tikzcd}
\]
It is worth checking the fibers directly.  Regard a point of \(Y\) as
the rank-at-most-one matrix
\[
\begin{pmatrix}
a&b\\
c&d
\end{pmatrix}
=
\begin{pmatrix}x_1\\x_3\end{pmatrix}
\begin{pmatrix}x_2&x_4\end{pmatrix}.
\]
Away from the zero matrix, both its column line and its row line are
uniquely determined.  Therefore \(\pi_+\) and \(\pi_-\) are
isomorphisms over \(Y\setminus\{0\}\).

Over the origin, the answer changes.  In the positive chamber,
\((x_1,x_3)\) cannot vanish, so the equations
\[
x_1x_2=x_1x_4=x_3x_2=x_3x_4=0
\]
force \(x_2=x_4=0\).  The remaining quotient is
\[
\pi_+^{-1}(0)
=
\{(x_1,x_3)\neq(0,0)\}/\Gm
\simeq\PP^1.
\]
Similarly,
\[
\pi_-^{-1}(0)
\simeq
\{(x_2,x_4)\neq(0,0)\}/\Gm
\simeq\PP^1.
\]
Each exceptional locus is a curve inside a threefold, hence has
codimension two.  Both contractions are small: no divisor disappears.
The birational map replaces the column-line \(\PP^1\) by the row-line
\(\PP^1\).

This is also exactly what happened to the square.  On one side we used
the diagonal \(13\); on the other we used the diagonal \(24\); on the
wall the diagonal disappeared.  The correspondence is
\[
\boxed{
\begin{array}{ccc}
\mathcal T_{13}
&\longleftrightarrow&
X_-
\\
\text{unsplit square}
&\longleftrightarrow&
Y
\\
\mathcal T_{24}
&\longleftrightarrow&
X_+.
\end{array}}
\]

Finally, the exceptional curve in either resolution is the zero
section of
\(\mathcal O_{\PP^1}(-1)^{\oplus2}\), so its normal bundle is
\[
N_{\PP^1/X_\pm}
\simeq
\mathcal O_{\PP^1}(-1)\oplus
\mathcal O_{\PP^1}(-1).
\]
Adjunction gives
\[
K_{X_\pm}|_{\PP^1}
\simeq
K_{\PP^1}\otimes
\det(N_{\PP^1/X_\pm})^{-1}
\simeq
\mathcal O_{\PP^1}.
\]
Thus \(K_{X_\pm}\cdot\PP^1=0\).  The elementary flip is crepant, so it
is in fact the standard threefold \emph{flop}.

\begin{agcheatsheet}
The terms \emph{flip} and \emph{flop} describe the same wall diagram
with different canonical behavior.
\[
\begin{array}{c|c}
\textbf{wall move}&\textbf{canonical degree on contracted curves}
\\
\hline
\text{flip, with a chosen direction}
&K\cdot C<0\text{ before and }K\cdot C'>0\text{ after}
\\
\text{flop}
&K\cdot C=K\cdot C'=0.
\end{array}
\]
In the toric setting, both are produced by replacing one simplicial
subdivision of a nonsimplicial cone by the subdivision on the other
side of its circuit relation.  The square relation has coefficient sum
zero, which is the combinatorial shadow of the crepant condition.
\end{agcheatsheet}

\subsection{Reading a wall crossing in four steps}

The examples suggest a workflow that applies well beyond them.
\begin{enumerate}
\item \textbf{Find the wall.}
      Locate a codimension-one face of a GIT chamber in character
      space.
\item \textbf{Compute the wall ring.}
      Determine the semi-invariants of degree \(d\theta_0\).  Form
      \(\Proj(R_{\theta_0})\), or take the invariant ring when the
      wall is \(\theta_0=0\).
\item \textbf{Find the exceptional locus.}
      Compare the chamber semistable locus with the wall semistable
      locus and determine which orbit closures acquire the same image.
\item \textbf{Classify the map.}
      A dimension drop is fiber type; a contracted divisor is
      divisorial; a codimension-at-least-two exceptional locus is
      small, and the neighboring chamber supplies the flip.
\end{enumerate}

\begin{pauseandthink}
In the \(\mathbb F_1\to\PP^2\) wall move, a ray disappears and the
exceptional curve is a divisor.  In the resolved-conifold wall move,
all four rays remain and only the triangulation changes.  Why do equal
ray sets but different maximal cones force the induced birational map
to be an isomorphism in codimension one?  Section~10.4 will identify
the contracted curves canonically by dualizing this nef wall to an
extremal ray of the Mori cone.
\end{pauseandthink}

\section{Extremal rays and contractions}
\label{sec:extremal-rays-contractions}
\index{Mori cone}
\index{nef cone}
\index{extremal ray}
\index{wall relation}
\index{extremal contraction}

Section~\ref{sec:crossing-wall} classified a wall morphism by looking
at its exceptional locus.  We now approach the same morphism from the
opposite direction.  Instead of starting with the map and asking what
it contracts, we start with the numerical classes of curves and ask
which map they determine.

The bridge is intersection theory.  A divisor on the boundary of the
nef cone has degree zero on some curves.  Those curve classes lie on a
face of the Mori cone.  In the toric setting, every nef divisor is
basepoint free, so its sections produce a morphism that contracts
exactly those curves.  Thus dual polyhedral cones turn a numerical
vanishing condition into an honest geometric contraction.

\begin{guidingquestion}
How does a facet of the nef cone identify the curves contracted by a
wall morphism, and how can those curve classes be computed directly
from the fan?
\end{guidingquestion}

\subsection{Divisors and curves modulo numerical equivalence}

Let \(X\) be a smooth projective variety.  Two Cartier divisors are
\emph{numerically equivalent} if they have the same intersection
number with every irreducible curve.  After tensoring with \(\RR\), the
space of numerical divisor classes is
\[
N^1(X)_\RR
=
\bigl(\Pic(X)/{\equiv_{\mathrm{num}}}\bigr)\otimes_\ZZ\RR.
\]
Similarly, two real one-cycles are numerically equivalent if every
Cartier divisor has the same degree on them.  The resulting space is
\[
N_1(X)_\RR
=
Z_1(X)_\RR/{\equiv_{\mathrm{num}}}.
\]
The definitions give a perfect pairing of finite-dimensional real
vector spaces
\[
\boxed{
N^1(X)_\RR\times N_1(X)_\RR
\longrightarrow\RR,
\qquad
([D],[C])\longmapsto D\cdot C.}
\]
Their common dimension is the \emph{Picard number}
\[
\rho(X)=\dim N^1(X)_\RR=\dim N_1(X)_\RR.
\]

Inside the curve space, form the cone generated by effective curves.

\begin{definition}[The Mori cone]
The \emph{Mori cone}, or \emph{cone of curves}, is
\[
\overline{\NE}(X)
=
\overline{
\Cone{[C]:C\subseteq X\text{ an irreducible curve}}}
\subseteq N_1(X)_\RR.
\]
A one-dimensional face of \(\overline{\NE}(X)\) is an
\emph{extremal ray}.
\end{definition}

The closure in this definition matters for general varieties.  For a
projective toric variety the cone will be rational polyhedral and
hence already closed.

A divisor class is nef precisely when it has nonnegative degree on
every curve.  Therefore
\[
\boxed{
\Nef(X)
=
\overline{\NE}(X)^\vee
=
\{\alpha\in N^1(X)_\RR:
  \alpha\cdot\gamma\geq0
  \text{ for every }\gamma\in\overline{\NE}(X)\}.}
\]
This equality is not an analogy: the nef and Mori cones are literal
dual cones under the intersection pairing.  In particular, a facet
of one is perpendicular to an extremal ray of the other.

\subsection{The toric cone theorem}

Let \(X_\Sigma\) be an \(n\)-dimensional smooth projective toric
variety.  A codimension-one cone
\[
\tau\in\Sigma(n-1)
\]
is traditionally called a \emph{wall of the fan}.  The orbit--cone
correspondence gives
\[
C_\tau:=V(\tau),
\qquad
\dim C_\tau=n-(n-1)=1,
\]
so every fan wall produces a complete torus-invariant curve.

These finitely many curves generate every numerical curve class that
is effective.

\begin{theorem}[The toric cone theorem]
\label{thm:toric-cone-theorem}
For a smooth projective toric variety \(X_\Sigma\),
\[
\boxed{
\overline{\NE}(X_\Sigma)
=
\sum_{\tau\in\Sigma(n-1)}
\RR_{\geq0}[V(\tau)].}
\]
Consequently \(\overline{\NE}(X_\Sigma)\) and
\(\Nef(X_\Sigma)\) are dual strongly convex rational polyhedral
cones of full dimension.  Every extremal ray is generated by the
class of a torus-invariant curve.
\end{theorem}

\noindent
\textit{Reference.}
The generation statement is Cox--Little--Schenck,
Theorem~6.3.20, and the projective duality statement is
Theorem~6.3.22.  Those results hold under weaker hypotheses than the
smooth projective case used here.

The theorem is the curve-side counterpart of a fact already used
throughout the notes: a toric divisor is nef if and only if its support
function is convex.  Convexity can be tested across the codimension-one
walls of the fan, and crossing such a wall measures the intersection
number with \(V(\tau)\).  Thus
\[
D\cdot V(\tau)\geq0
\quad\text{for every }\tau\in\Sigma(n-1)
\quad\Longleftrightarrow\quad
D\text{ is nef}.
\]
Dualizing these finitely many inequalities gives the displayed
description of the Mori cone.

Two cautions are useful.
First, every extremal ray is generated by an invariant curve, but an
arbitrary invariant curve need not generate an extremal ray.  Several
invariant curves may also be numerically equivalent and hence
determine the same ray.  Second, the word ``wall'' now occurs in three
closely related places.

\begin{agcheatsheet}
Keep the three walls distinct.
\[
\begin{array}{c|c|c}
\textbf{wall}&\textbf{where it lives}&\textbf{what it determines}
\\
\hline
\tau\in\Sigma(n-1)
&\text{the fan in }N_\RR
&\text{the invariant curve }V(\tau)
\\
\text{facet of }\Nef(X)
&N^1(X)_\RR
&\text{an extremal ray of }\overline{\NE}(X)
\\
\text{wall of }\Sigma_{\mathrm{GKZ}}
&\text{character space}
&\text{a limiting GIT quotient.}
\end{array}
\]
For a Cox chamber with all rays present, the last two walls are the
same wall under the identification of that chamber with the nef cone.
The first wall supplies an invariant curve representing the dual
extremal ray.
\end{agcheatsheet}

\subsection{Wall relations compute curve classes}

The class \([V(\tau)]\) can be read from one local relation among ray
generators.  Let
\[
\tau=\Cone{u_2,\dots,u_n}
\]
be a wall of a smooth complete fan, and write the two adjacent maximal
cones as
\[
\sigma=\Cone{u_1,u_2,\dots,u_n},
\qquad
\sigma'=\Cone{u_2,\dots,u_n,u_{n+1}}.
\]
Smoothness gives a unique integral relation of the form
\[
\boxed{
u_1+b_2u_2+\cdots+b_nu_n+u_{n+1}=0,}
\qquad b_i\in\ZZ.
\]
This is the \emph{wall relation} of \(\tau\).

\begin{proposition}[The wall relation as an intersection vector]
\label{prop:wall-relation-intersections}
Let \(C_\tau=V(\tau)\), and let \(D_i\) be the invariant divisor
corresponding to the ray through \(u_i\).  Then
\[
D_1\cdot C_\tau=D_{n+1}\cdot C_\tau=1,
\qquad
D_i\cdot C_\tau=b_i
\quad(2\leq i\leq n),
\]
and every other invariant divisor has intersection zero with
\(C_\tau\).  Hence, for
\(D=\sum_\rho a_\rho D_\rho\),
\[
\boxed{
D\cdot C_\tau
=
a_1+b_2a_2+\cdots+b_na_n+a_{n+1}.}
\]
\end{proposition}

\begin{proof}[Why this formula is natural]
Intersecting a principal divisor with \(C_\tau\) gives zero.  Thus for
every \(m\in M\),
\[
0
=
\operatorname{div}(\chi^m)\cdot C_\tau
=
\sum_\rho
\langle m,u_\rho\rangle(D_\rho\cdot C_\tau).
\]
Equivalently,
\[
\sum_\rho(D_\rho\cdot C_\tau)u_\rho=0
\]
in \(N\).  The orbit--cone correspondence shows that only the rays in
\(\sigma\cup\sigma'\) can contribute.  Moreover, \(D_1\) and
\(D_{n+1}\) meet \(C_\tau\) transversely at its two torus-fixed
points, so both intersections are one.  The unique wall relation then
forces all the remaining coefficients.
\end{proof}

\noindent
\textit{Reference.}
This is the smooth case of Cox--Little--Schenck,
Proposition~6.4.4; their formula includes the cone multiplicities
needed in the simplicial singular case.  The interpretation of a
numerical curve class as a relation among ray generators is
Equation~(15.1.4).

For a surface, this recovers a familiar calculation.  If the rays are
cyclically ordered and
\[
u_{i-1}+u_{i+1}=a_i u_i,
\]
then the wall relation for the invariant curve \(D_i\) is
\[
u_{i-1}-a_i u_i+u_{i+1}=0.
\]
Proposition~\ref{prop:wall-relation-intersections} therefore gives
\[
D_i^2=-a_i.
\]
The surface self-intersection formula from Chapter~8 is thus the
two-dimensional instance of the general wall-relation calculation.

\subsection{A nef facet produces an extremal contraction}

Let \(R\subseteq\overline{\NE}(X)\) be an extremal ray.  Its dual
facet is
\[
F_R
=
\{\alpha\in\Nef(X):\alpha\cdot\gamma=0
  \text{ for every }\gamma\in R\}.
\]
Choose a rational divisor class \([D]\) in the relative interior of
\(F_R\), and replace it by a positive integral multiple.  Since \(D\)
is nef and toric, it is basepoint free.  Its section ring therefore
defines a morphism
\[
\varphi_D:X
\longrightarrow
X_D
:=
\Proj\!\left(
\bigoplus_{k\geq0}H^0(X,\mathcal O_X(kD))
\right).
\]

\begin{theorem}[Toric extremal contraction]
\label{thm:toric-extremal-contraction}
Let \(X_\Sigma\) be a smooth projective toric variety and let
\(R\subseteq\overline{\NE}(X_\Sigma)\) be an extremal ray.  There is a
projective toric morphism with connected fibers
\[
\varphi_R:X_\Sigma\longrightarrow X_0
\]
such that, for every irreducible curve \(C\subseteq X_\Sigma\),
\[
\boxed{
\varphi_R(C)\text{ is a point}
\quad\Longleftrightarrow\quad
[C]\in R.}
\]
The morphism may be obtained from any integral nef class in the
relative interior of the dual facet \(F_R\).
\end{theorem}

\noindent
\textit{Reference.}
This is the projective smooth case of Cox--Little--Schenck,
Proposition~15.4.1, which treats simplicial semiprojective toric
varieties.  The equivalence ``nef = basepoint free'' is
Theorem~6.3.12.

Here is the mechanism.  The morphism associated to a basepoint-free
toric divisor factors \(D\), up to linear equivalence, as
\[
D\sim\varphi_D^*A
\]
for an ample divisor \(A\) on the image.  Therefore
\[
D\cdot C=0
\quad\Longleftrightarrow\quad
\varphi_D(C)\text{ is a point}.
\]
Because \([D]\) lies in the relative interior of \(F_R\), the only
effective classes on which it vanishes are those in \(R\).  This proves
the numerical characterization of the contracted curves.

The same statement can be seen directly in the fan and polytope.  The
normal fan of \(P_D\) is obtained by joining adjacent maximal cones
whose Cartier data agree.  Across a fan wall
\(\tau=\sigma\cap\sigma'\),
\[
m_\sigma=m_{\sigma'}
\quad\Longleftrightarrow\quad
D\cdot V(\tau)=0.
\]
Joining those cones makes \(V(\tau)\) invisible to the resulting
toric morphism.  This is exactly the fan map constructed by the
extremal contraction theorem.

More generally, if \([D]\) lies in the relative interior of a
lower-dimensional face \(F\subseteq\Nef(X)\), then it contracts the
entire dual face
\[
F^\perp\cap\overline{\NE}(X).
\]
An \emph{elementary} contraction is the special case in which \(F\)
is a facet and the dual face is one extremal ray.

\subsection{Surface laboratory: the dual cones of
\texorpdfstring{\(\mathbb F_1\)}{F1}}

We can now recover both morphisms from Section~10.3 by one two-by-two
intersection calculation.  Write
\[
\mathbb F_1=\Bl_p(\PP^2),
\qquad
E^2=-1,
\qquad
H^2=1,
\qquad
H\cdot E=0,
\]
and let
\[
F=H-E
\]
be the class of a fiber of the ruling.  To keep the two numerical
spaces visually distinct, write
\[
e=[E]\in N_1(\mathbb F_1)_\RR,
\qquad
f=[F]\in N_1(\mathbb F_1)_\RR
\]
for the curve classes, while \(H,F\in N^1(\mathbb F_1)_\RR\) denote
divisor classes.  Their pairing is
\[
\begin{array}{c|cc}
&e=[E]&f=[F]
\\
\hline
H&0&1
\\
F=H-E&1&0.
\end{array}
\]

Every effective curve class is a nonnegative combination of \(e\) and
\(f\), so
\[
\boxed{
\overline{\NE}(\mathbb F_1)
=
\Cone{e,f}.}
\]
Indeed, if \(D=\alpha H+\beta F\), then the table gives
\[
D\cdot E=\beta,
\qquad
D\cdot F=\alpha.
\]
Hence
\[
\boxed{
\Nef(\mathbb F_1)
=
\Cone{H,F},}
\]
and \(D\) is ample exactly when \(\alpha,\beta>0\).

The two boundary rays of the nef cone are perpendicular to the two
extremal rays of the Mori cone:
\[
\begin{array}{c|c|c|c}
\textbf{nef facet}
&\textbf{vanishing}
&\textbf{dual extremal ray}
&\textbf{contraction}
\\
\hline
\RR_{\geq0}H
&H\cdot E=0
&\RR_{\geq0}e
&\mathbb F_1\to\PP^2
\\
\RR_{\geq0}F
&F\cdot F=0
&\RR_{\geq0}f
&\mathbb F_1\to\PP^1.
\end{array}
\]
Thus the blowdown and the ruling are not two unrelated constructions.
They are the extremal contractions belonging to the two edges of one
nef cone.

This also clarifies their different geometric types.  The ray
\(\RR_{\geq0}e\) contains the class of a single irreducible curve,
the exceptional section \(E\); contracting it is divisorial.  The ray
\(\RR_{\geq0}f\) contains the class of every fiber in a moving family
covering \(\mathbb F_1\); contracting all of them lowers dimension and
is of fiber type.

\begin{surfacelaboratory}
The same computation works for every Hirzebruch surface.  If \(S\) is
the section with \(S^2=-r\) and \(F\) is a fiber, then
\[
\overline{\NE}(\mathbb F_r)=\Cone{[S],[F]},
\qquad
\Nef(\mathbb F_r)=\Cone{F,S+rF}.
\]
Indeed, for \(D=aS+bF\),
\[
D\cdot F=a,
\qquad
D\cdot S=b-ar.
\]
The two inequalities \(a\geq0\) and \(b-ar\geq0\) cut out the nef
cone.  When \(r=1\), the class \(S+F=H\) gives the blowdown.  When
\(r>1\), the same nef boundary contracts the negative section to a
singular torus-fixed point rather than to a smooth point.
\end{surfacelaboratory}

\subsection{The conifold relation revisited}

The square example from Section~10.3 shows the same dictionary in
dimension three.  On the side using the diagonal \(24\), the two
maximal cones meet along
\[
\tau=\Cone{\nu_2,\nu_4},
\]
and the wall relation is
\[
\nu_1-\nu_2+\nu_3-\nu_4=0.
\]
For the exceptional curve \(C=V(\tau)\simeq\PP^1\),
Proposition~\ref{prop:wall-relation-intersections} gives
\[
(D_1\cdot C,D_2\cdot C,D_3\cdot C,D_4\cdot C)
=(1,-1,1,-1).
\]
The curve space is one-dimensional, so
\[
\overline{\NE}(X_+)=\RR_{\geq0}[C].
\]
Moving the linearization to the origin reaches the unique nef wall and
contracts precisely this ray, producing the quadric cone
\(Y=V(ad-bc)\).

The wall relation also detects crepancy without using the normal bundle.
Since
\[
K_{X_+}=-D_1-D_2-D_3-D_4,
\]
we obtain
\[
K_{X_+}\cdot C
=
-(1-1+1-1)
=0.
\]
On the other resolution the relation, and hence the curve class, is
reversed.  The canonical degree is still zero.  This is the numerical
reason that the wall crossing is a flop.

\subsection{Nef chambers inside the secondary fan}

We can now close the circle begun in Section~10.2.  For a simplicial
semiprojective toric quotient in which every Cox coordinate remains
present, the corresponding GKZ cone is naturally identified with the
nef cone of the quotient:
\[
\boxed{
\Gamma_{\Sigma,\varnothing}
\simeq
\Nef(X_\Sigma).}
\]
This is Cox--Little--Schenck, Proposition~15.1.3.  Under this
identification, the three constructions from this chapter line up:
\[
\begin{array}{c|c|c}
\textbf{secondary fan}
&\textbf{divisors and curves}
&\textbf{geometry}
\\
\hline
\text{chamber for }X_\Sigma
&\Nef(X_\Sigma)
&\text{one polarized toric model}
\\
\text{facet of the chamber}
&R^\perp\cap\Nef(X_\Sigma)
&\text{extremal contraction of }R
\\
\text{neighboring chamber}
&\text{nef cone of the next model}
&\text{flip or another birational model.}
\end{array}
\]

This also explains the subtlety noticed for \(\mathbb F_1\) in
Section~10.2.  Its chamber \(\Cone{F,H}\) is exactly
\(\Nef(\mathbb F_1)\).  Across the \(H\)-wall, the coordinate \(v_0\)
is forced to be nonzero and its ray disappears.  The adjacent
two-dimensional GKZ chamber therefore contains an extra GIT direction;
after removing that redundant direction, the quotient \(\PP^2\) has
the expected one-dimensional nef cone.

\begin{pauseandthink}
Suppose \(K_X\) is not nef.  By duality there is an extremal ray
\(R\subseteq\overline{\NE}(X)\) with
\[
K_X\cdot R<0.
\]
Theorem~\ref{thm:toric-extremal-contraction} now supplies a canonical
next geometric move: contract \(R\), and flip if the contraction is
small.  Repeating this choice is the minimal model program.  What does
this algorithm look like as a path through the secondary fan, and
where can it end for a projective toric variety?  Section~10.5 will
assemble the final picture.
\end{pauseandthink}

\section{The toric minimal model program}
\label{sec:toric-mmp}
\index{minimal model program!toric}
\index{Mori fiber space}
\index{termination!toric flips}
\index{canonical divisor!MMP direction}

We now have all of the ingredients promised at the beginning of the
chapter.  A Cox character chooses a GIT quotient, the secondary fan
organizes those quotients into chambers, a chamber is the nef cone of
one model, and a facet of that nef cone is dual to an extremal ray of
curves.  The minimal model program adds one final instruction:
\emph{among all extremal rays, move only in a direction on which the
canonical divisor is negative.}

We begin with a smooth projective toric variety, but we must immediately
allow a slightly larger class of spaces.  A divisorial contraction or a
flip can create singularities even when the starting variety is smooth.
The natural setting is therefore a projective \(\QQ\)-factorial toric
variety with terminal singularities.  For a toric variety,
\(\QQ\)-factoriality is equivalent to the fan being simplicial.

\begin{guidingquestion}
How does the sign of \(K_X\cdot R\) orient the wall crossings in the
secondary fan, why must the resulting process stop, and what geometric
object is left at the end?
\end{guidingquestion}

\subsection{The algorithm: contract, flip, or stop}

Let \(X\) be a projective \(\QQ\)-factorial toric variety with terminal
singularities.  If \(K_X\) is not nef, the toric cone theorem supplies
an extremal ray
\[
R\subseteq\overline{\NE}(X)
\qquad\text{with}\qquad
K_X\cdot R<0.
\]
The contraction theorem then gives a projective toric morphism with
connected fibers
\[
\varphi_R:X\longrightarrow Y
\]
that contracts exactly the curves whose classes lie in \(R\).  There
are three possibilities.

\begin{enumerate}
\item If \(\dim Y<\dim X\), then \(\varphi_R\) is of \emph{fiber type}.
The process stops.

\item If \(\varphi_R\) is birational and contracts a divisor, replace
\(X\) by \(Y\) and continue.  This is a \emph{divisorial step}.

\item If \(\varphi_R\) is birational but contracts no divisor, then it
is small.  Replace \(X\) by the other side of the toric flip
\[
\begin{tikzcd}[column sep=huge]
X \arrow[dr,"\varphi_R"']
  \arrow[rr,dashed,"\psi"]
&& X^+ \arrow[dl,"\varphi_R^+"]
\\
&Y.&
\end{tikzcd}
\]
and continue on \(X^+\).
\end{enumerate}

In the last case the contracted curve class changes sign.  More
precisely, if \(R^+\) is the extremal ray of \(X^+\) contracted by
\(\varphi_R^+\), then
\[
K_X\cdot R<0
\qquad\Longrightarrow\qquad
K_{X^+}\cdot R^+>0.
\]
Thus a flip does not merely replace one triangulation by another: it
turns a \(K\)-negative direction into a \(K\)-positive direction.

\begin{definition}[Mori fiber space]
A \emph{Mori fiber space} is a contraction
\[
f:X\longrightarrow Z
\]
with connected fibers such that
\[
\dim Z<\dim X,
\qquad
\rho(X)-\rho(Z)=1,
\qquad
-K_X\text{ is }f\text{-ample}.
\]
For a \(K_X\)-negative extremal contraction of fiber type, these
conditions are automatic.
\end{definition}

If instead the process reaches a model \(X'\) for which \(K_{X'}\) is
nef, then \(X'\) is called a \emph{minimal model}.  The general MMP has
these two possible outcomes: a minimal model or a Mori fiber space.
For projective toric varieties, however, there is a stronger conclusion.

\begin{theorem}[The projective toric \(K\)-MMP]
\label{thm:projective-toric-mmp}
Starting from a positive-dimensional projective \(\QQ\)-factorial toric
variety with terminal singularities, any sequence of \(K\)-negative
divisorial contractions and flips terminates after finitely many steps
in a Mori fiber space.
\end{theorem}

\begin{proof}[Why the two assertions hold]
First, a positive-dimensional projective toric variety cannot have nef
canonical divisor.  If \(A\) is ample and
\(K_X=-\sum_{\rho}D_\rho\), then
\[
K_X\cdot A^{n-1}
=
-\sum_\rho D_\rho\cdot A^{n-1}<0,
\]
because every invariant prime divisor has positive intersection with
\(A^{n-1}\).  Hence the algorithm cannot end with \(K_X\) nef; when it
stops, it must stop at a fiber-type contraction.

It remains to know that it stops.  Every divisorial step decreases the
Picard number by one, so only finitely many divisorial steps can occur.
Between them, flips preserve the set of rays of the fan, and there are
only finitely many fans built from a fixed finite ray set.  Finiteness
alone does not exclude a cycle, but \(K\)-negativity does.  On a common
toric refinement of a \(K\)-negative flip, the pullbacks of the two
canonical divisors differ by a nonzero effective exceptional divisor.
Going around a cycle would therefore say that one fixed pullback of
\(K\) equals itself plus a nonzero effective divisor, which is
impossible.  Thus there is no infinite sequence of flips.
\end{proof}

\noindent
\textit{Reference.}
The toric procedure is Cox--Little--Schenck,
Procedure~15.5.5.  Proposition~15.5.3 shows that \(K\)-negative
divisorial contractions and flips preserve terminal singularities, and
Theorem~15.5.10 proves termination of toric flips.  Their results are
stated in the more general semiprojective and relative setting.

\subsection{The MMP as a directed path through the secondary fan}

The algorithm has a particularly concrete GIT interpretation.  Let
\(A\) be an ample divisor on the current model and consider the segment
in divisor-class space
\[
\theta(t)=(1-t)[A]+t[K_X],
\qquad 0\leq t\leq1.
\]
It begins in the interior of \(\Nef(X)\).  Suppose it first reaches the
boundary at \(t=t_0\), on the facet dual to an extremal ray \(R\).  Then
\[
0=\theta(t_0)\cdot R
=(1-t_0)A\cdot R+t_0K_X\cdot R,
\]
and therefore
\[
\boxed{
K_X\cdot R
=
-\frac{1-t_0}{t_0}\,A\cdot R<0.}
\]
So moving from an ample class toward the canonical class automatically
selects a \(K\)-negative wall.

What happens at that wall is exactly the trichotomy from
Section~\ref{sec:crossing-wall}.

\begin{agcheatsheet}
\textbf{The toric MMP in the secondary fan.}

\medskip

\begin{center}
\small
\begin{tabularx}{\linewidth}{@{}
  >{\raggedright\arraybackslash}p{0.23\linewidth}
  >{\raggedright\arraybackslash}p{0.28\linewidth}
  >{\raggedright\arraybackslash}X@{}}
\toprule
\textbf{wall reached}
&\textbf{fan / GIT change}
&\textbf{MMP instruction}
\\
\midrule
flipping wall
&replace one circuit triangulation by the other
&cross to the adjacent nef chamber and continue
\\
divisorial wall
&delete the ray of the exceptional divisor
&lower the Picard number, then rebuild the chamber picture
\\
boundary of the effective region
&the quotient becomes degenerate in one direction
&take the fiber-type contraction and stop.
\\
\bottomrule
\end{tabularx}
\end{center}
\end{agcheatsheet}

The word \emph{directed} is essential.  VGIT allows us to cross a wall
in either direction.  The \(K\)-MMP allows only the direction for which
the curves being removed have negative canonical degree.  It is this
numerical sign that turns the secondary fan from a catalogue of models
into an algorithm.

\subsection{A complete run on \texorpdfstring{\(\mathbb F_1\)}{F1}}

Return one last time to
\[
\mathbb F_1=\Bl_p(\PP^2).
\]
Use the curve classes
\[
e=[E],
\qquad
f=[F],
\qquad
\overline{\NE}(\mathbb F_1)=\Cone{e,f}
\]
from Section~\ref{sec:extremal-rays-contractions}.  Since
\[
K_{\mathbb F_1}=-3H+E=-2E-3F,
\]
we obtain
\[
K_{\mathbb F_1}\cdot E=-1,
\qquad
K_{\mathbb F_1}\cdot F=-2.
\]
Both extremal rays are \(K\)-negative, so the first MMP step is not
unique.

\[
\begin{tikzcd}[column sep=huge,row sep=large]
&\mathbb F_1
  \arrow[dl,"\text{contract }E"']
  \arrow[dr,"\text{ruling}"]
&
\\
\PP^2 \arrow[d,"\text{to a point}"']
&&\PP^1
\\
\{\mathrm{pt}\}&&
\end{tikzcd}
\]

On the left branch, contracting \(e\) is the divisorial blowdown
\[
\mathbb F_1\longrightarrow\PP^2.
\]
The Mori cone of \(\PP^2\) has one ray, generated by the class \(\ell\)
of a line, and
\[
K_{\PP^2}\cdot\ell=-3.
\]
Its extremal contraction is \(\PP^2\to\{\mathrm{pt}\}\), a Mori fiber
space.

On the right branch, choosing \(f\) gives the ruling immediately:
\[
\mathbb F_1\longrightarrow\PP^1.
\]
This is already a Mori fiber space, since its relative Picard number is
one and
\[
-K_{\mathbb F_1}\cdot F=2>0.
\]
Thus one starting variety can have more than one valid MMP.  The two
runs even end over bases of different dimensions.  The algorithm is
finite and tightly constrained, but it is not canonical until the
extremal ray is chosen.

\begin{surfacelaboratory}
The same calculation for \(\mathbb F_r\) exposes an important warning.
If \(S^2=-r\) and \(F\) is a fiber, then
\[
K_{\mathbb F_r}=-2S-(r+2)F,
\]
so
\[
K_{\mathbb F_r}\cdot S=r-2,
\qquad
K_{\mathbb F_r}\cdot F=-2.
\]
Hence
\[
\begin{array}{c|c|c|c}
r&K\cdot S&K\cdot F&\text{allowed first \(K\)-MMP step}
\\
\hline
0&-2&-2&\text{either ruling}
\\
1&-1&-2&\text{blow down \(S\), or take the ruling}
\\
2&0&-2&\text{the ruling only}
\\
r\geq3&r-2>0&-2&\text{the ruling only.}
\end{array}
\]
For \(r\geq2\), the negative section still spans an extremal ray of the
Mori cone, but that ray is not \(K\)-negative.  In particular, on
\(\mathbb F_2\) its contraction is \(K\)-trivial, not a step of the
ordinary \(K\)-MMP.  A curve having negative self-intersection is not
the same thing as its ray being \(K\)-negative.
\end{surfacelaboratory}

\subsection{Flips are directed; flops are not MMP steps}

The sign test can also be read directly from a wall relation.  Suppose
an invariant curve \(C\) is represented by
\[
\sum_i b_i u_i=0
\]
and the coefficients \(b_i\) are its intersections with the
corresponding invariant divisors.  Since
\(K_X=-\sum_iD_i\),
\[
\boxed{
K_X\cdot C=-\sum_i b_i.}
\]
Thus the circuit itself tells us whether the associated wall crossing
is permitted by the \(K\)-MMP:
\[
\sum_i b_i>0
\quad\Longleftrightarrow\quad
K_X\cdot C<0.
\]
If the contraction is small, crossing the wall reverses the circuit
class.  The new exceptional curves consequently have positive
canonical degree, which records the direction of the flip.

Now compare the conifold wall from Sections~\ref{sec:crossing-wall}
and~\ref{sec:extremal-rays-contractions}.  Its relation is
\[
\nu_1-\nu_2+\nu_3-\nu_4=0.
\]
The coefficient sum is zero, and hence
\[
K_X\cdot C=-(1-1+1-1)=0.
\]
The two small resolutions are related by a \emph{flop}, not by a
\(K\)-negative flip.  VGIT sees the two neighboring chambers perfectly
well, but the ordinary \(K\)-MMP does not prefer either direction.
Flops belong to the comparison of minimal models, where the canonical
class is numerically unchanged; they are adjacent to the MMP story,
but they are not descent steps in the ordinary \(K\)-MMP.

\subsection{The chapter's final dictionary}

We can now read one geometric move in every language developed in the
notes.  Begin with a toric variety \(X_\Sigma\) and a \(K\)-negative
extremal ray \(R\).  Then
\[
\boxed{
\begin{aligned}
\text{fan:}\quad
&\text{a wall relation represents }R,
\\
\text{intersection theory:}\quad
&K_X\cdot R<0,
\\
\text{nef--Mori duality:}\quad
&R^\perp\cap\Nef(X)\text{ is a facet},
\\
\text{Cox/GIT:}\quad
&\text{the corresponding character reaches a wall},
\\
\text{secondary fan:}\quad
&\text{cross, delete a ray, or reach the boundary},
\\
\text{birational geometry:}\quad
&\text{flip, contract a divisor, or form a Mori fiber space.}
\end{aligned}}
\]

This is the promised finale.  The Cox quotient from Chapter~6 did not
merely construct one toric variety.  Once its character is allowed to
move, it organizes a finite birational world:
\[
\boxed{
\text{Cox quotient}
\longrightarrow
\text{secondary fan}
\longrightarrow
\text{wall crossing}
\longrightarrow
\text{toric MMP}.}
\]

\begin{pauseandthink}
The two runs from \(\mathbb F_1\) show exactly where choice enters the
MMP: both rays \(e\) and \(f\) are \(K\)-negative, and the canonical
class does not choose between them.  What the theory guarantees is not
a unique path, but a legal move at every stage and finite termination.

To test whether you now own the full dictionary, start with a wall
relation, compute its canonical degree, identify the dual nef facet,
and predict whether the corresponding GIT wall gives a divisorial
contraction, a flip, a flop, or a Mori fiber space.  The concluding
exercises will ask you to carry out precisely this chain without being
told which language to begin in.
\end{pauseandthink}

\clearpage
\section*{What this chapter has accomplished}
\addcontentsline{toc}{section}{What this chapter has accomplished}
\markboth{What this chapter has accomplished}{What this chapter has accomplished}

The chapter began by returning to the Cox construction and changing
only one ingredient.  We fixed an affine space with a diagonal torus
action, but allowed the character
\[
\theta\in\Hom(G,\Gm)_\RR
\]
to move.  The basic calculation was the weight-cone criterion
\[
\boxed{
z\text{ is }\theta\text{-semistable}
\quad\Longleftrightarrow\quad
\theta\in
\Cone{w_i:z_i\neq0}.}
\]
It turns semistability into convex geometry: a point survives exactly
when the weights of its nonzero coordinates can support the chosen
character.  For an ample Cox character this recovers the usual Cox
open set, while other characters may produce different toric
quotients from the same polynomial ring and the same group action.

The regions on which the semistable locus remains constant assemble
into the secondary, or GKZ, fan.  Its support is the effective weight
cone, its maximal cones are GIT chambers, and its faces record limiting
quotients.  Gale duality translates this chamber structure into
regular subdivisions of a ray or point configuration.  In the
affine-hyperplane setting, the same incidence data are encoded dually
by the secondary polytope:
\[
\boxed{
\begin{gathered}
\text{GIT chamber}
\longleftrightarrow
\text{regular triangulation}
\longleftrightarrow
\text{vertex of }P_{\mathrm{GKZ}},
\\
\text{GIT wall}
\longleftrightarrow
\text{coarser regular subdivision}
\longleftrightarrow
\text{edge or higher face of }P_{\mathrm{GKZ}}.
\end{gathered}}
\]

Crossing a wall put geometry back into this combinatorics.  A
fiber-type wall lowers the dimension of the quotient; a divisorial
wall removes a ray and contracts an invariant divisor; a small wall
keeps the rays and changes how they are assembled into cones.  In the
last case the two neighboring models are related by a flip, or by a
flop when the contracted curves have canonical degree zero.  The two
boundary rays of \(\Nef(\mathbb F_1)\) displayed the fiber-type and
divisorial possibilities, while the two resolutions of the conifold
displayed the first genuinely three-dimensional flop.

The numerical meaning of a wall came from nef--Mori duality:
\[
\boxed{
\Nef(X)=\overline{\NE}(X)^\vee.}
\]
For a projective toric variety, invariant curves generate the Mori
cone.  A fan wall gives such a curve, its wall relation gives the
intersection vector of that curve, and the perpendicular nef facet
gives the morphism that contracts it.  Thus the three different walls
used in the chapter---a codimension-one cone of the fan, a facet of the
nef cone, and a wall of the GKZ fan---are three appearances of one
numerical contraction.

Finally, the canonical class oriented the chamber picture.  If a wall
relation represents an invariant curve \(C\), then
\[
K_X\cdot C=-\sum_i(D_i\cdot C).
\]
The toric \(K\)-MMP repeatedly chooses a ray on which this number is
negative and then contracts it, flips it, or stops at a Mori fiber
space.  Divisorial steps decrease the Picard number, flips cannot cycle,
and the process terminates.  Moreover, a positive-dimensional
projective toric variety cannot have nef canonical divisor, so the
absolute toric \(K\)-MMP ends with a Mori fiber space.  In the language
of this chapter, it is a directed path through finitely many chambers.

The exercises first test each part of this dictionary separately.  The
challenge then asks you to reconstruct an entire toric flip from a
single one-dimensional weight matrix.

\begin{exercises}[Comprehension exercises]
\markboth{Comprehension exercises}{Comprehension exercises}

\item Let \(\Gm\) act on \(\AAA^3\) with coordinates \((x_0,x_1,y)\)
and weights
\[
1,\qquad1,\qquad-1.
\]
Use the weight-cone criterion to determine the semistable loci for
characters \(\theta>0\) and \(\theta<0\).  Show that the corresponding
quotients are
\[
\operatorname{Tot}\bigl(\mathcal O_{\PP^1}(-1)\bigr)
\qquad\text{and}\qquad
\AAA^2,
\]
respectively.  Compute the invariant ring at \(\theta=0\), identify
the wall quotient with \(\AAA^2\), and explain why the wall morphism
from the positive chamber contracts the zero section of
\(\mathcal O_{\PP^1}(-1)\).

\item Return to the Cox weights of \(\mathbb F_1\),
\[
F,\qquad S,\qquad F,\qquad F+S.
\]
Starting only from these four weights, use the weight-cone criterion
to recover the two chambers
\[
a>b>0,
\qquad
b>a>0
\]
for \(\theta=aF+bS\).  Determine which coordinate patterns are
unstable in each chamber.  Then explain why the first quotient is
\(\mathbb F_1\), the second is \(\PP^2\), and the common wall gives
the blowdown of the exceptional curve.

\item For the square configuration
\[
(0,0),\quad(1,0),\quad(1,1),\quad(0,1),
\]
write down its two triangulations.  Compute the GKZ vector of each
triangulation using normalized lattice area, and verify that their
convex hull is a segment.  Show directly that the normal fan of this
segment has two chambers separated by one wall.  Which subdivision of
the square corresponds to the wall itself?

\item Consider the two circuit relations
\[
u_1+u_3=u_2+u_4
\qquad\text{and}\qquad
u_1+u_2+u_3=u_4+u_5.
\]
Regard their coefficients as intersection vectors of contracted
curves.  Compute \(K_X\cdot C\) in each case.  Explain why the first
relation is compatible with a flop, whereas the second has a preferred
\(K\)-MMP direction.  What additional geometric information is needed
to decide whether the second contraction is divisorial or small?

\item Let \(X=\PP^1\times\PP^1\), with ruling curve classes
\(f_1,f_2\) and divisor classes \(H_1,H_2\).  Use
\[
H_i\cdot f_j=\delta_{ij}
\]
to prove
\[
\overline{\NE}(X)=\Cone{f_1,f_2},
\qquad
\Nef(X)=\Cone{H_1,H_2}.
\]
Identify the extremal contraction determined by each nef boundary
ray.  Since \(K_X=-2H_1-2H_2\), show that either contraction is a
valid first and final step of the \(K\)-MMP.

\item Use the standard fan of \(\PP^3\), whose rays are generated by
\[
e_1,\quad e_2,\quad e_3,\quad-e_1-e_2-e_3.
\]
Take the fan wall \(\tau=\Cone{e_1,e_2}\).  Write the wall relation
using the two adjacent rays, and use it to compute the intersections
of the invariant curve \(C_\tau\) with all four invariant divisors.
Deduce
\[
-K_{\PP^3}\cdot C_\tau=4.
\]
Identify \(C_\tau\) geometrically and reconcile the answer with
\(-K_{\PP^3}=4H\).

\item On the Hirzebruch surface \(\mathbb F_r\), let \(S^2=-r\) and
let \(F\) be a ruling fiber.  Starting from
\[
K_{\mathbb F_r}=-2S-(r+2)F,
\]
determine which extremal rays are \(K\)-negative for every
\(r\geq0\).  Run every possible ordinary \(K\)-MMP beginning at
\(\mathbb F_0\), \(\mathbb F_1\), and \(\mathbb F_2\).  Explain why
the contraction of the negative section on \(\mathbb F_2\) belongs to
the VGIT picture but not to the ordinary \(K\)-MMP.

\item Let \(X\) be a positive-dimensional projective toric variety
and let \(A\) be ample.  Use
\[
K_X=-\sum_{\rho}D_\rho
\]
to show
\[
K_X\cdot A^{n-1}<0.
\]
Why does this rule out \(K_X\) being nef?  Explain why an absolute
toric \(K\)-MMP can therefore terminate only with a fiber-type
contraction.  Finally, explain which part of the argument would have
to change for a relative MMP over a nontrivial base.

\end{exercises}

\clearpage
\section*{Challenge exercise}
\addcontentsline{toc}{section}{Challenge exercise}
\markboth{Challenge exercise}{Challenge exercise}

\noindent
\textbf{Build a toric flip from one line of weights.}

\medskip

Let \(G=\Gm\) act on
\[
\AAA^5=\Spec\CC[x_0,x_1,x_2,y_0,y_1]
\]
by
\[
\lambda\cdot(x_0,x_1,x_2,y_0,y_1)
=
(\lambda x_0,\lambda x_1,\lambda x_2,
 \lambda^{-1}y_0,\lambda^{-1}y_1).
\]
Thus the weight matrix is the single row
\[
Q=\begin{pmatrix}1&1&1&-1&-1\end{pmatrix}.
\]
There are two nonzero chambers in
\(\Hom(G,\Gm)_\RR\simeq\RR\).  The goal is to reconstruct the two
quotients, their common wall contraction, and the direction of the
resulting flip using only this row of weights.

\begin{enumerate}[label=\textbf{Hint \arabic*.},wide=0pt]

\item \textbf{Find the two semistable loci.}

Apply the weight-cone criterion.  Prove that for \(\theta>0\),
\[
U_+
=
\{(x,y):(x_0,x_1,x_2)\neq(0,0,0)\},
\]
whereas for \(\theta<0\),
\[
U_-
=
\{(x,y):(y_0,y_1)\neq(0,0)\}.
\]
Draw the one-dimensional GKZ fan: two rays separated by the wall
\(\theta=0\).  Which coordinate subspace is removed in each chamber?

\item \textbf{Identify the chamber quotients.}

For \(\theta>0\), the nonzero vector
\(x=(x_0,x_1,x_2)\) determines a point of \(\PP^2\).  Check from the
\(G\)-action that each \(y_j\) is a fiber coordinate of
\(\mathcal O_{\PP^2}(-1)\).  Deduce
\[
\boxed{
X_+:=U_+/G
\simeq
\operatorname{Tot}\!\left(
\mathcal O_{\PP^2}(-1)^{\oplus2}\right).}
\]
Reverse the roles of \(x\) and \(y\) to prove
\[
\boxed{
X_-:=U_-/G
\simeq
\operatorname{Tot}\!\left(
\mathcal O_{\PP^1}(-1)^{\oplus3}\right).}
\]
In particular, both quotients are smooth fourfolds.

\item \textbf{Compute the wall quotient.}

At \(\theta=0\), the quotient is affine:
\[
X_0=\Spec\CC[x_0,x_1,x_2,y_0,y_1]^G.
\]
Show that the invariant ring is generated by the six quadratic
monomials
\[
z_{ij}=x_i y_j,
\qquad
0\leq i\leq2,
\quad
0\leq j\leq1.
\]
Arrange them in a \(3\times2\) matrix \(Z=(z_{ij})\).  Prove that the
relations are the \(2\times2\) minors of \(Z\), and conclude
\[
\boxed{
X_0
=
\{Z\in\operatorname{Mat}_{3\times2}(\CC):
  \rank(Z)\leq1\}.}
\]
Equivalently, \(X_0\) is the affine cone over the Segre embedding
\(\PP^2\times\PP^1\subset\PP^5\).

\item \textbf{Find the exceptional loci.}

The invariant functions define morphisms
\[
\begin{tikzcd}[column sep=huge]
X_+ \arrow[dr,"\pi_+"']
&&
X_- \arrow[dl,"\pi_-"]
\\
&X_0.&
\end{tikzcd}
\]
Show that \(\pi_+\) contracts the zero section
\(E_+\simeq\PP^2\) to the vertex of the cone and that \(\pi_-\)
contracts the zero section \(E_-\simeq\PP^1\).  Prove that each map is
an isomorphism away from its zero section.  Since
\[
\codim_{X_+}(E_+)=2,
\qquad
\codim_{X_-}(E_-)=3,
\]
deduce that both contractions are small.  Thus the induced map
\[
X_+\dashrightarrow X_-
\]
is an elementary toric flip in one of its two orientations.

\item \textbf{Let the canonical class choose the orientation.}

For a vector bundle \(p:E\to B\), use the tangent sequence of its
total space to prove
\[
K_{\operatorname{Tot}(E)}
\simeq
p^*\!\left(K_B\otimes(\det E)^{-1}\right).
\]
Apply this formula to the two chamber quotients.  On \(X_+\), show
\[
K_{X_+}|_{E_+}
\simeq
\mathcal O_{\PP^2}(-1),
\]
so \(K_{X_+}\cdot L=-1\) for a line
\(L\subset E_+\).  On \(X_-\), show
\[
K_{X_-}|_{E_-}
\simeq
\mathcal O_{\PP^1}(1),
\]
so \(K_{X_-}\cdot E_-=1\).  Conclude that
\[
\boxed{
X_+\dashrightarrow X_-
\text{ is directed from the \(K\)-negative side to the
\(K\)-positive side}.}
\]

\item \textbf{Recover the sign from the circuit.}

Let \(u_0,\ldots,u_4\) be the images of the standard basis of
\(\ZZ^5\) in the quotient lattice determined by \(Q\).  Show that the
unique circuit relation is
\[
u_0+u_1+u_2-u_3-u_4=0.
\]
Its coefficient sum is \(1\).  Use the wall-relation formula from
Section~10.4 to recover
\[
K_{X_+}\cdot L=-1.
\]
Explain how the two subdivisions of this circuit keep the same rays
but assemble them into different maximal cones, exactly as required
for a small wall crossing.

\item \textbf{Compare a flip with a flop.}

Replace the weights by
\[
(1,1,-1,-1).
\]
Repeat the canonical-class calculation, or simply use the coefficient
sum of the circuit relation, to show that the two chamber quotients
are \(K\)-trivial over their common affine quotient.  Recover the
resolved-conifold flop from Section~10.3.  Which single numerical
feature of the weight row distinguishes the flop from the flip above?

\item \textbf{Read the whole construction in six languages.}

Summarize the example by completing the following chain without
looking back at the preceding hints:
\[
\boxed{
\begin{gathered}
\text{weight row}
\longrightarrow
\text{semistable loci}
\longrightarrow
\text{two GKZ chambers}
\\[-1mm]
\longrightarrow
\text{two circuit subdivisions}
\longrightarrow
\text{small contractions}
\longrightarrow
\text{\(K\)-directed flip}.
\end{gathered}}
\]
At which arrow is convex geometry used?  Where does invariant theory
enter?  Where is the orbit--cone dictionary hiding?  Which calculation
turns an undirected VGIT wall crossing into an MMP step?

\end{enumerate}

\begin{threedictionaries}
The entire flip can be recovered from the row
\(Q=(1,1,1,-1,-1)\).

\medskip

\begin{center}
\small
\begin{tabularx}{\linewidth}{@{}
  >{\raggedright\arraybackslash}X
  >{\raggedright\arraybackslash}X
  >{\raggedright\arraybackslash}X@{}}
\toprule
\textbf{Weights / fan}
&\textbf{GIT quotient}
&\textbf{Birational geometry}
\\
\midrule
positive chamber
&\(\operatorname{Tot}(\mathcal O_{\PP^2}(-1)^{\oplus2})\)
&small \(K\)-negative contraction
\\
wall
&cone over \(\PP^2\times\PP^1\)
&common contraction target
\\
negative chamber
&\(\operatorname{Tot}(\mathcal O_{\PP^1}(-1)^{\oplus3})\)
&flipped, \(K\)-positive model
\\
coefficient sum \(1\)
&unbalanced circuit
&flip rather than flop
\\
\bottomrule
\end{tabularx}
\end{center}
\end{threedictionaries}

This example contains the whole chapter in miniature.  A grading
produces a torus action; a character selects an open set; the chambers
organize the possible quotients; a circuit describes the wall; and
one canonical intersection number directs the birational move.

\clearpage
\section*{A final perspective}
\addcontentsline{toc}{section}{A final perspective}
\markboth{A final perspective}{A final perspective}

These notes began with a strongly convex rational polyhedral cone.
From its lattice points we built a semigroup algebra and an affine
toric variety.  Gluing cones into a fan produced global varieties;
support functions and polytopes produced divisors, line bundles, and
projective embeddings; Cox coordinates replaced gluing by one
multigraded quotient.  The compact torus then turned the polytope into
an orbit space, the fan became a cohomology ring, and localization and
Riemann--Roch turned characteristic classes into lattice-point counts.
In this final chapter, varying the Cox character made the same data
move through an entire birational world.

The progression can be compressed into one last diagram:
\[
\boxed{
\begin{gathered}
\text{cone and semigroup}
\longrightarrow
\text{affine toric variety},
\\
\text{fan and support function}
\longrightarrow
\text{global variety and line bundle},
\\
\text{polytope and Cox grading}
\longrightarrow
\text{sections, topology, and quotients},
\\
\text{moving character}
\longrightarrow
\text{secondary fan and toric MMP}.
\end{gathered}}
\]

What makes toric geometry unusually coherent is not merely that its
objects admit pictures.  It is that the same finite combinatorial data
can be read, without loss, in several mathematical languages.  A ray
may be a divisor, a Cox variable, or a facet normal.  A wall may be an
invariant curve, a nef equation, or a change of quotient.  A lattice
point may be a character, a monomial section, or a contribution to an
index.  Moving between these descriptions is the central skill the
book has tried to develop.

Toric varieties are therefore both a testing ground and a source of
models.  They make singularities, resolutions, positivity,
cohomology, localization, and the minimal model program explicit
enough to compute, while retaining the structure of the general
theories they illuminate.  The combinatorics does not replace the
geometry.  It reveals how the geometry is organized.

\clearpage
\section*{Notes and references}
\addcontentsline{toc}{section}{Notes and references}
\markboth{Notes and references}{Notes and references}

The GIT construction used in Section~10.1 is developed in
Cox--Little--Schenck, \S14.1.  Proposition~14.1.9 identifies the
semistable locus of an ample Cox character with the usual Cox open
set.  Definition~14.1.11 and Proposition~14.1.12 give the
section-ring construction of the GIT quotient; the weight-cone test
used throughout the chapter is its direct monomial specialization,
proved in Section~10.1.

The GKZ decomposition of the effective cone is treated in
Cox--Little--Schenck, \S14.4; Theorem~14.4.7 states that the GKZ cones
form a fan and that the quotient is constant on the relative interior
of each cone.  Regular triangulations and the secondary polytope are
developed in their \S15.2.  We used Definition~15.2.8 and
Proposition~15.2.9 for the passage between chambers and regular
triangulations, and Proposition~15.2.11 together with
Corollary~15.2.12 for the secondary polytope and its vertices.

Our wall-crossing discussion follows Cox--Little--Schenck, \S15.3 and
\S15.4.  Theorem~15.3.13 describes the birational maps attached to a
flipping wall, Proposition~15.4.5 gives the fiber-type, divisorial, and
small trichotomy, and Proposition~15.4.1 constructs the contraction
dual to a nef facet.  The resolved conifold is the basic balanced
circuit: its two triangulations give small crepant resolutions, so the
wall crossing is a flop.

The numerical input comes from the toric intersection theory of
Cox--Little--Schenck, Chapter~6.  Theorems~6.3.20 and~6.3.22 give the
generation of the Mori cone by invariant curves and its duality with
the nef cone.  Proposition~6.4.4 computes the intersection vector of
an invariant curve from its wall relation, and Theorem~6.3.12 supplies
the crucial toric fact that nef divisors are basepoint free.

The toric MMP appears in Cox--Little--Schenck, \S15.5.
Procedure~15.5.5 gives the algorithm, Proposition~15.5.3 records
preservation of terminal singularities under the relevant divisorial
contractions and flips, and Theorem~15.5.10 proves termination of
toric flips.  Our treatment concentrated on the absolute projective
case; the same combinatorics also works in relative and
semiprojective settings.

The final challenge is the standard unbalanced circuit with weight row
\((1,1,1,-1,-1)\).  Its wall quotient is the determinantal cone of
rank-at-most-one \(3\times2\) matrices.  The two small resolutions have
exceptional loci \(\PP^2\) and \(\PP^1\), and the unequal numbers of
positive and negative weights make the canonical degree nonzero.  The
balanced row \((1,1,-1,-1)\) instead gives the conifold flop.

The notes have emphasized a single habit of thought: whenever a toric
calculation appears in one language, look for its counterparts in the
lattice, fan, polytope, divisor, Cox, topological, and birational
languages.  That dictionary is more durable than any isolated formula,
and it is the real conclusion of the book.
